BOOK ONE

The Doctrine of Being

WITH WHAT MUST THE BEGINNING OF SCIENCE BE MADE?

It is only in recent times that there has been a new awareness of
the difficulty of finding a beginning in philosophy,
and the reason for this difficulty,
and so also the possibility of resolving it,
have been discussed in a variety of ways.
The beginning of philosophy must be either
something mediated or something immediate,
and it is easy to show that it can be neither the one nor the other;
so either way of beginning runs into contradiction.

The principle of a philosophy also expresses a beginning, of course,
but not so much a subjective as an objective one,
the beginning of all things.
The principle is a somehow determinate content:
“water,” “the one,” “nous,” “idea,” or “substance,” “monad,” etc.
or, if it designates the nature of cognition
and is therefore meant simply as a criterion
rather than an objective determination,
as “thinking,” “intuition,” “sensation,” “I,” even “subjectivity,”
then here too the interest still lies in the content determination.
The beginning as such, on the other hand, as something subjective
in the sense that it is an accidental way of introducing the exposition,
is left unconsidered, a matter of indifference,
and consequently also the need to ask
with what a beginning should be made remains
of no importance in face of the need for the principle
in which alone the interest of the fact seems to lie,
the interest as to what is the truth,
the absolute ground of everything.

But the modern perplexity about a beginning proceeds from
a further need which escapes those who are either busy demonstrating
their principle dogmatically or skeptically looking for
a subjective criterion against dogmatic philosophizing,
and is outright denied by those who begin, like a shot from a pistol,
from their inner revelation, from faith, intellectual intuition, etc.
and who would be exempt from method and logic.
If earlier abstract thought is at first interested only
in the principle as content, but is driven as philosophical culture
advances to the other side to pay attention to the conduct of
the cognitive process, then the subjective activity has
also been grasped as an essential moment of objective truth,
and with this there comes the need to unite
the method with the content, the form with the principle.
Thus the principle ought to be also the beginning,
and that which has priority for thinking ought to be also
the first in the process of thinking.

Here we only have to consider how the logical beginning appears.
The two sides from which it can be taken have already been named,
namely either by way of mediation as result,
or immediately as beginning proper.
This is not the place to discuss the question
apparently so important to present-day culture,
whether the knowledge of truth is
an immediate awareness that begins absolutely, a faith,
or rather a mediated knowledge.
In so far as the issue allows passing treatment,
this has already been done elsewhere
(in my Encyclopedia of the Philosophical Sciences, 3rd edn, in the
Prefatory Concept, §§21ff.).
Here we may quote from it only this, that
there is nothing in heaven or nature or spirit or anywhere else
that does not contain just as much immediacy as mediation,
so that both these determinations prove to be unseparated and inseparable
and the opposition between them nothing real.
As for a scientific discussion, a case in point is
every logical proposition in which we find
the determinations of immediacy and mediacy
and where there is also entailed, therefore,
a discussion of their opposition and their truth.
This opposition, when connected to thinking,
to knowledge, to cognition,
assumes the more concrete shape of immediate or mediated knowledge,
and it is then up to the science of logic to consider
the nature of cognition in general,
while the more concrete forms of
the same cognition fall within the scope of
the science of spirit and the phenomenology of spirit.
But to want to clarify the nature of cognition
prior to science is to demand that it should be discussed outside science,
and outside science this cannot be done,
at least not in the scientific manner
which alone is the issue here.

A beginning is logical in that it is to be made in
the element of a free, self-contained thought, in pure knowledge;
it is thereby mediated, for pure knowledge is
the ultimate and absolute truth of consciousness.
We said in the Introduction that the Phenomenology of Spirit is
the science of consciousness, its exposition;
that consciousness has the concept of science,
that is, pure knowledge, for its result.
To this extent, logic has for its presupposition
the science of spirit in its appearance,
a science which contains the necessity,
and therefore demonstrates the truth,
of the stand-point which is pure knowledge and of its mediation.
In this science of spirit in its appearance
the beginning is made from empirical, sensuous consciousness,
and it is this consciousness which is immediate knowledge in the strict sense;
there, in this science, is where its nature is discussed.
Any other consciousness, such as faith in divine truths, inner experience,
knowledge through inner revelation, etc., proves upon cursory reflection
to be very ill-suited as an instance of immediate knowledge.
In the said treatise, immediate consciousness is also
that which in the science comes first and immediately
and is therefore a presupposition;
but in logic the presupposition is what has proved itself to be
the result of that preceding consideration,
namely the idea as pure knowledge.
Logic is the pure science, that is,
pure knowledge in the full compass of its development.
But in that result the idea has the determination of
a certainty that has become truth;
it is a certainty which, on the one hand, no longer stands
over and against a subject matter confronting it externally
but has interiorized it, is knowingly aware
that the subject matter is itself;
and, on the other hand, has relinquished
any knowledge of itself that would oppose it to objectivity
and would reduce the latter to a nothing;
it has externalized this subjectivity
and is at one with its externalization.

Now starting with this determination of pure knowledge,
all that we have to do to ensure that the beginning will remain
immanent to the science of this knowledge is to consider,
or rather, setting aside every reflection, simply to take up,
what is there before us.
Pure knowledge, thus withdrawn into this unity,
has sublated every reference to an other and to mediation;
it is without distinctions and as thus distinctionless
it ceases to be knowledge;
what we have before us is only simple immediacy.

Simple immediacy is itself an expression of reflection;
it refers to the distinction from what is mediated.
The true expression of this simple immediacy is therefore pure being.
Just as pure knowledge should mean nothing but knowledge as such,
so also pure being should mean nothing but being in general;
being, and nothing else, without further determination and filling.

Being is what makes the beginning here; it is presented indeed as origi-
nating through mediation, but a mediation which at the same time sublates
itself, and the presupposition is of a pure knowledge which is the result
of finite knowledge, of consciousness. But if no presupposition is to be
made, if the beginning is itself to be taken immediately, then the only
determination of this beginning is that it is to be the beginning of logic, of
thought as such. There is only present the resolve, which can also be viewed
as arbitrary, of considering thinking as such. The beginning must then be
absolute or, what means the same here, must be an abstract beginning; and
so there is nothing that it may presuppose, must not be mediated by anything
or have a ground, ought to be rather itself the ground of the entire science.
It must therefore be simply an immediacy, or rather only immediacy itself.
Just as it cannot have any determination with respect to an other, so too it
cannot have any within; it cannot have any content, for any content would
entail distinction and the reference of distinct moments to each other, and
hence a mediation. The beginning is therefore pure being.

After this simple exposition of what alone first belongs to this simplest of
all simples, the logical beginning, we may add the following further reflec-
tions which should not serve, however, as elucidation and confirmation of
the exposition – this is complete by itself – but are rather occasioned by
notions and reflections which may come our way beforehand and yet, like
all other prejudices that antedate the science of logic, must be disposed
of within the science itself and are therefore to be patiently deferred until
then.

The insight that absolute truth must be a result, and conversely, that a
result presupposes a first truth which, because it is first, objectively con-
sidered is not necessary and from the subjective side is not known – this
insight has recently given rise to the thought that philosophy can begin only
with something which is hypothetically and problematically true, and that
at first, therefore, philosophizing can be only a quest. This is a view that
Reinhold has repeatedly urged in the later stages of his philosophizing, 4
and which must be given credit for being motivated by a genuine interest
in the speculative nature of philosophical beginning. A critical examination
of this view will also be an occasion for introducing a preliminary under-
standing of what progression in logic generally means, for the view has
direct implications for the nature of this advance. Indeed, as portrayed by
it, progression in philosophy would be rather a retrogression and a ground-
ing, only by virtue of which it then follows as result that that, with which
the beginning was made, was not just an arbitrary assumption but was in
fact the truth, and the first truth at that.

It must be admitted that it is an essential consideration
(one which will be found elaborated again within the logic itself)
that progression is a retreat to the ground,
to the origin and the truth on which
that with which the beginning was made,
and from which it is in fact produced, depends.
Thus consciousness, on its forward path from the immediacy with which it began,
is led back to the absolute knowledge which is its innermost truth.
This truth, the ground, is then also that from which the original first proceeds,
the same first which at the beginning came on the scene as something immediate.
It is most of all in this way that absolute spirit
(which is revealed as the concrete and supreme truth of all being)
comes to be known, as at the end of the development it freely externalizes itself,
letting itself go into the shape of an immediate being,
resolving itself into the creation of a world
which contains all that fell within the development
preceding that result and which,
through this reversal of position with its beginning,
is converted into something dependent on the result as principle.
Essential to science is not so much that
a pure immediacy should be the beginning,
but that the whole of science is in itself a circle
in which the first becomes also the last,
and the last also the first.

Conversely, it follows that it is just as necessary to consider as result that
into which the movement returns as to its ground. In this respect, the first
is just as much the ground, and the last a derivative; since the movement
makes its start from the first and by correct inferences arrives at the last as the
ground, this last is result. Further, the advance from that which constitutes
the beginning is to be considered only as one more determination of the
same advance, so that this beginning remains as the underlying ground of
all that follows without vanishing from it. The advance does not consist in
the derivation of an other, or in the transition to a truly other: inasmuch
as there is a transition, it is equally sublated again. Thus the beginning
of philosophy is the ever present and self-preserving foundation of all
subsequent developments, remaining everywhere immanent in its further
determinations.

In this advance the beginning thus loses the one-sidedness that it has
when determined simply as something immediate and abstract;
it becomes mediated, and the line of scientific forward movement consequently turns
into a circle.
It also follows that what constitutes the beginning,
because it is something still undeveloped and empty of content,
is not yet truly known at that beginning,
and that only science, and science fully developed, is
the completed cognition of it,
replete with content and finally truly grounded.

But for this reason, because it is as absolute ground that the result finally
emerges, the progression of this cognition is not anything provisory, still
problematic and hypothetical, but must be determined through the nature
of the matter at issue and of the content itself. Nor is the said beginning an
arbitrary and only temporary assumption, 5 or something which seems to
be an arbitrary and tentative presupposition but of which it is subsequently
shown that to make it the starting point was indeed the right thing to do;
this is not as when we are instructed to make certain constructions in order
to aid the proof of a geometrical theorem, and only in retrospect, in the
course of the proof, does it become apparent that we did well to draw
precisely these lines and then, in the proof itself, to begin by comparing
them or the enclosed angles – though the line-drawing or the comparing
themselves escape conceptual comprehension.

So we have just given, right within science itself, the reason why in pure
science the beginning is made with pure being. This pure being is the unity
into which pure knowledge returns, or if this knowledge, as form, is itself
still to be kept distinct from its unity, then pure being is also its content.
It is in this respect that this pure being, this absolute immediate, is just as
absolutely mediated. However, just because it is here as the beginning, it is
just as essential that it should be taken in the one-sidedness of being purely
immediate. If it were not this pure indeterminacy, if it were determined, it
would be taken as something mediated, would already be carried further
than itself: a determinate something has the character of an other with
respect to a first. It thus lies in the nature of a beginning itself that it should
be being and nothing else. There is no need, therefore, of other preparations
to enter philosophy, no need of further reflections or access points.

Nor can we derive a more specific determination or a more positive content
for the beginning of philosophy from the fact that it is such a beginning.
For here, at the beginning, where the fact itself is not yet at hand, philosophy
is an empty word, a received and yet unjustified notion. Pure knowledge
yields only this negative determination, namely that the beginning ought
to be abstract. If pure being is taken as the content of pure knowledge,
then the latter must step back from its content, allowing it free play and
without determining it further. – Or again, inasmuch as pure being is
to be considered as the unity into which knowledge has collapsed when
at the highest point of union with its objectification, knowledge has then
disappeared into this unity, leaving behind no distinction from it and hence
no determination for it. – Nor is there anything else present, any content
whatever, that could be used to make a more determinate beginning with

But, it may be said, the determination of being assumed so far as the
beginning can also be let go, so that the only requirement would be that a
pure beginning should be made. Nothing would then be at hand except the
beginning itself, and we must see what this would be. – This position could
be suggested also for the benefit of those who are either not comfortable,
for whatever reason, with beginning with being and even less with the
transition into nothing that follows from being, or who simply do not
know how else to make a beginning in a science except by presupposing
a representation which is subsequently analyzed, the result of the analysis
then yielding the first determinate concept in the science. If we also want
to test this strategy, we must relinquish every particular object that we may
intend, since the beginning, as the beginning of thought, is meant to be
entirely abstract, entirely general, all form with no content; we must have
nothing, therefore, except the representation of a mere beginning as such.
We have, therefore, only to see what there is in this representation.

As yet there is nothing, and something is supposed to become. The
beginning is not pure nothing but a nothing, rather, from which something
is to proceed; also being, therefore, is already contained in the beginning.
Therefore, the beginning contains both, being and nothing; it is the unity
of being and nothing, or is non-being which is at the same time being, and
being which is at the same time non-being.

Further, being and nothing are present in the beginning as distinguished;
for the beginning points to something other – it is a non-being which refers
to an other; that which begins, as yet is not; it only reaches out to being.
The being contained in the beginning is such, therefore, that it distances
itself from non-being or sublates it as something which is opposed to it.
But further, that which begins already is, but is also just as much not yet.
The opposites, being and non-being, are therefore in immediate union in
it; or the beginning is their undifferentiated unity.

An analysis of the beginning would thus yield the concept of the unity of
being and non-being – or, in a more reflected form, the concept of the unity
of differentiated and undifferentiated being – or of the identity of identity
and non-identity. 7 This concept could be regarded as the first, purest, that
is, most abstract, definition of the absolute – as it would indeed be if the
issue were just the form of definitions and the name of the absolute. In this
sense, just as such an abstract concept would be the first definition of the
absolute, so all further determinations and developments would be only
more determinate and richer definitions of it. But let those who are not
satisfied with being as the beginning, since being passes over into nothing
and what emerges is the unity of the two – let them consider what is more
likely to satisfy them: this beginning that begins with the representation
of the beginning and an analysis of it (an analysis that is indeed correct yet
equally leads to the unity of being and non-being) or a beginning which
makes being the beginning.

But, regarding this strategy, there is still a further observation to be made.
The said analysis presupposes that the representation of the beginning is
known; its strategy follows the example of other sciences. These presuppose
their object and presume that everyone has the same representation of it and
will find in it roughly the same determinations which they have collected
here or there, through analysis, comparison, and sundry argumentation,
and they then offer as its representations. But that which constitutes the
absolute beginning must likewise be something otherwise known; now,
if it is something concrete and hence in itself variously determined, then
this connectedness which it is in itself is presupposed as a known; the con-
nectedness is thereby adduced as something immediate, which however it is
not; for it is connectedness only as a connection of distinct elements and
therefore contains mediation within itself. Further, the accidentality and
the arbitrariness of the analysis and the specific mode of determination
affect the concrete internally. Which determinations are elicited depends
on what each individual happens to discover in his immediate accidental
representation. The connection contained within a concrete something,
within a synthetic unity, is necessary only in so far as it is not found already
given but is produced rather by the spontaneous return of the moments
back into this unity, a movement which is the opposite of the analytical
procedure that occurs rather within the subject and is external to the fact
itself.

Here we then have the precise reason why that with which the begin-
ning is to be made cannot be anything concrete, anything containing a
connection within its self. It is because, as such, it would presuppose within
itself a process of mediation and the transition from a first to an other, of
which process the concrete something, now become a simple, would be
the result. But the beginning ought not itself to be already a first and an
other, for anything which is in itself a first and an other implies that an
advance has already been made. Consequently, that which constitutes the
beginning, the beginning itself, is to be taken as something unanalyzable,
taken in its simple, unfilled immediacy; and therefore as being, as complete
emptiness.

If, impatient with this talk of an abstract beginning, one should say that
the beginning is to be made, not with the beginning, but directly with the
fact itself, well then, this subject matter is nothing else than that empty
being. For what this subject matter is, that is precisely what ought to result
only in the course of the science, what the latter cannot presuppose to
know in advance.

On any other form otherwise assumed in an effort to have a beginning
other than empty being, that beginning would still suffer from the same
defects. Let those who are still dissatisfied with this beginning take upon
themselves the challenge of beginning in some other way and yet avoiding
such defects.

But we cannot leave entirely unmentioned a more original beginning
to philosophy which has recently gained notoriety, the beginning with the
“I.” It derived from both the reflection that all that follows from the first
truth must be deduced from it, and the need that this first truth should be
something with which one is already acquainted, and even more than just
acquainted, something of which one is immediately certain. This proposed
beginning is not, as such, an accidental representation, or one which might
be one thing to one subject and something else to another. For the “I,” this
immediate consciousness of the self, appears from the start to be both itself
an immediate something and something with which we are acquainted in a
much deeper sense than with any other representation; true, anything else
known belongs to this “I,” but it belongs to it as a content which remains
distinct from it and is therefore accidental; the “I,” by contrast, is the simple
certainty of its self. But the “I” is, as such, at the same time also a concrete,
or rather, the “I” is the most concrete of all things – the consciousness of
itself as an infinitely manifold world. Before the “I” can be the beginning
and foundation of philosophy, this concreteness must be excised, and this
is the absolute act by virtue of which the “I” purifies itself and makes its
entrance into consciousness as abstract “I.” But this pure “I” is now not
immediate, is not the familiar, ordinary “I” of our consciousness to which
everyone immediately links science. Truly, that act of excision would be
none other than the elevation to the standpoint of pure knowledge in
which the distinction between subject and object has disappeared. But
as thus immediately demanded, this elevation is a subjective postulate;
before it proves itself as a valid demand, the progression of the concrete “I”
from immediate consciousness to pure knowledge must be demonstratively
exhibited within the “I” itself, through its own necessity. Without this
objective movement, pure knowledge, also when defined as intellectual
intuition, appears as an arbitrary standpoint, itself one of those empirical
states of consciousness for which everything depends on whether someone,
though not necessarily somebody else, discovers it within himself or is able
to produce it there. But inasmuch as this pure “I” must be essential, pure
knowledge – and pure knowledge is however one which is only posited in
individual consciousness through an absolute act of self-elevation, is not
present in it immediately – we lose the very advantage which was to derive
from this beginning of philosophy, namely that it is something with which
everyone is well acquainted, something which everyone finds within himself
and to which he can attach further reflection; that pure “I,” on the contrary,
in its abstract, essential nature, is to ordinary consciousness an unknown,
something that the latter does not find within itself. What comes with it is
rather the disadvantage of the illusion that we are speaking of something
supposedly very familiar, the “I” of empirical self-consciousness, whereas
at issue is in fact something far removed from the latter. Determining
pure knowledge as “I” acts as a continuing reminder of the subjective “I”
whose limitations should rather be forgotten; it leads to the belief that
the propositions and relations which result from the further development
of the “I” occur within ordinary consciousness and can be found pre-
given there, indeed that the whole issue is about this consciousness. This
mistake, far from bringing clarity, produces instead an even more glaring
and bewildering confusion; among the public at large, it has occasioned
the crudest of misunderstandings.

Further, as regards the subjective determinateness of the “I” in general,
pure knowledge does remove from it the restriction that it has when under-
stood as standing in unsurmountable opposition to an object. But for this
reason it would be at least superfluous still to hold on to this subjective
attitude by determining pure knowledge as “I.” For this determination not
only carries with it that troublesome duality of subject and object; on closer
examination, it also remains a subjective “I.” The actual development of
the science that proceeds from the “I” shows that in the course of it the
object has and retains the self-perpetuating determination of an other with
respect to the “I”; that therefore the “I” from which the start was made
does not have the pure knowledge that has truly overcome the opposition
of consciousness, but is rather still entangled in appearance.
In this connection, there is the further essential observation to be made
that, although the “I” might well be determined to be in itself pure knowl-
edge or intellectual intuition and declared to be the beginning, in science
we are not concerned with what is present in itself or as something inner, bu
with the external existence 9 rather of what in thought is inner and with the
determinateness which this inner assumes in that existence. But whatever
externalization there might be of 10 intellectual intuition at the beginning of
science, or – if the subject matter of science is called the eternal, the divine,
the absolute – of the eternal or absolute, this cannot be anything else than a
first, immediate, simple determination. Whatever richer name be given to
it than is expressed by mere being, the only legitimate consideration is how
such an absolute enters into discursive 11 knowledge and the enunciation of
this knowledge. Intellectual intuition might well be the violent rejection
of mediation and of demonstrative, external reflection. However, anything
which it says over and above simple immediacy would be something con-
crete, and this concrete would contain a diversity of determinations in it.
But, as already remarked, the enunciation and exposition of this concrete
something is a process of mediation which starts with one of the determina-
tions and proceeds to another, even though this other returns to the first –
and this is a movement which, moreover, is not allowed to be arbitrary
or assertoric. Consequently, that from which the beginning is made in any
such exposition is not something itself concrete but only the simple imme-
diacy from which the movement proceeds. Besides, what is lacking if we
make something concrete the beginning is the demonstration which the
combination of the determinations contained in it requires.

Therefore, if in the expression of the absolute, or the eternal, or God
(and God would have the perfectly undisputed right
that the beginning be made with him),
if in the intuition or the thought of them,
there is more than there is in pure being,
then this more should first emerge in a knowledge
which is discursive and not figurative;
as rich as what is implicitly contained in knowledge may be,
the determination that first emerges in it is something simple,
for it is only in the immediate that
no advance is yet made from one thing to an other.
Consequently, whatever in the richer representations
of the absolute or God might be said or implied over and above being,
all this is at the beginning only an empty word and only being;
this simple determination which has no further meaning
besides, this empty something, is as such, therefore,
the beginning of philosophy.

This insight is itself so simple that this beginning is
as beginning in no need of any preparation or further introduction,
and the only possible purpose of this preliminary disquisition regarding it
was not to lead up to it but to dispense rather with all preliminaries.

GENERAL DIVISION OF BEING

Being is determined, first, as against another in general;
secondly, it is internally self-determining;
thirdly, as this preliminary division is cast off,
it is the abstract indeterminateness and immediacy
in which it must be the beginning.

According to the first determination,
being partitions itself off from essence,
for further on in its development it proves to be
in its totality only one sphere of the concept,
and to this sphere as moment it opposes another sphere.
According to the second, it is the sphere
within which fall the determinations
and the entire movement of its reflection.

In this, being will posit itself in three determinations:

I. as determinateness; as such, quality;

II. as sublated determinateness; magnitude, quantity;

III. as qualitatively determined quantity; measure.

This division, as was generally remarked of such divisions in the
Introduction, is here a preliminary statement;
its determinations must first arise from the movement of being itself,
and receive their definitions and justification by virtue of it.
As regards the divergence of this division from the usual listing
of the categories,  namely quantity, quality, relation and modality,
(for Kant, incidentally, these are supposed to be only classifications
of his categories, but are in fact themselves categories,
only more abstract ones; about this, there is nothing to remark here,
since the entire listing will diverge from the usual ordering and meaning
of the categories at every point.)

This only can perhaps be remarked, that the determination of quantity
is ordinarily listed ahead of quality and as a rule
this is done for no given reason.
It has already been shown that the beginning is made with being as such,
and hence with qualitative being.
It is clear from a comparison of quality with quantity that
the former is by nature first.
For quantity is quality which has already become negative;
magnitude is the determinateness which, no longer one with being but
already distinguished from it, is the sublated quality
that has become indifferent.
It includes the alterability of being without altering the fact itself,
namely being, of which it is the determination;
qualitative determinateness is on the contrary one with its being,
it neither transcends it nor stays within it
but is its immediate restrictedness.
Hence quality, as the determinateness which is immediate,
is the first and it is with it that the beginning is to be made.

Measure is a relation, not relation in general
but specifically of quality and quantity to each other;
the categories dealt with by Kant under relation
will come up elsewhere in their proper place.
Measure, if one so wishes, can be considered also a modality;
but since with Kant modality is no longer supposed to make up
a determination of content,
but only concerns
the reference of the content to thought, to the subjective,
the result is a totally heterogeneous reference
that does not belong here.

The third determination of being falls within
the section Quality inasmuch as being, as abstract immediacy,
reduces itself to one single determinateness
as against its other determinacies inside its sphere.

SECTION I

Determinateness (Quality)

Being is the indeterminate immediate;
it is free of determinateness with respect to essence,
just as it is still free of any determinateness
that it can receive within itself.
This reflectionless being is being
as it immediately is only within.

Since it is immediate, it is being without quality;
but the character of indeterminateness attaches to it in itself
only in opposition to what is determinate or qualitative.
Determinate being thus comes to stand over and against being in general;
with that, however, the very indeterminateness
of being constitutes its quality.
It will therefore be shown that the first being is
in itself determinate, and therefore secondly,
that it passes over into existence, is existence;
that this latter, however, as finite being, sublates itself
and passes over into the infinite reference of being to itself;
it passes over, thirdly, into being-for-itself.

CHAPTER 1

Being

A. BEING

Being, pure being, without further determination.
In its indeterminate immediacy it is equal only to itself
and also not unequal with respect to another;
it has no difference within it, nor any outwardly.
If any determination or content were posited in it as distinct,
or if it were posited by this determination or content
as distinct from an other,
it would thereby fail to hold fast to its purity.
It is pure indeterminateness and emptiness.
There is nothing to be intuited in it,
if one can speak here of intuiting;
or, it is only this pure empty intuiting itself.
Just as little is anything to be thought in it,
or, it is equally only this empty thinking.
Being, the indeterminate immediate is in fact nothing,
and neither more nor less than nothing.

B. NOTHING

Nothing, pure nothingness;
it is simple equality with itself,
complete emptiness,
complete absence of determination and content;
lack of all distinction within.
In so far as mention can be made here of
intuiting and thinking,
it makes a difference whether something or nothing is
being intuited or thought.
To intuit or to think nothing has therefore a meaning;
the two are distinguished and so nothing is (concretely exists)
in our intuiting or thinking;
or rather it is the empty intuiting and thinking itself,
like pure being.
Nothing is therefore the same determination
or rather absence of determination,
and thus altogether the same as what pure being is.

C. BECOMING

1. Unity of being and nothing

Pure being and pure nothing are therefore the same.
The truth is neither being nor nothing,
but rather that being has passed over into nothing and
nothing into being; “has passed over,” not passes over.
But the truth is just as much that
they are not without distinction;
it is rather that they are not the same,
that they are absolutely distinct
yet equally unseparated and inseparable,
and that each immediately vanishes in its opposite.
Their truth is therefore this movement of
the immediate vanishing of the one into the other:
becoming, a movement in which the two are distinguished,
but by a distinction which has just as immediately
dissolved itself.

2. The moments of becoming

Becoming is the unseparatedness of being and nothing,
not the unity that abstracts from being and nothing;
as the unity of being and nothing
it is rather this determinate unity,
or one in which being and nothing equally are.
However, inasmuch as being and nothing are
each unseparated from its other, each is not.
In this unity, therefore, they are,
but as vanishing, only as sublated.
They sink from their initially represented self-subsistence
into moments which are still distinguished
but at the same time sublated.

Grasped as thus distinguished,
each is in their distinguishedness
a unity with the other.
Becoming thus contains being and nothing as two such unities,
each of which is itself unity of being and nothing;
the one is being as immediate and as reference to nothing;
the other is nothing as immediate and as reference to being;
in these unities the determinations are of unequal value.

Becoming is in this way doubly determined.
In one determination, nothing is the immediate,
that is, the determination begins with nothing
and this refers to being;
that is to say, it passes over into it.
In the other determination, being is the immediate,
that is, the determination begins with being
and this passes over into nothing:
coming-to-be and ceasing-to-be.

Both are the same, becoming,
and even as directions that are so different
they interpenetrate and paralyze each other.
The one is ceasing-to-be;
being passes over into nothing,
but nothing is just as much the opposite of itself,
the passing-over into being, coming-to-be.
This coming-to-be is the other direction;
nothing goes over into being,
but being equally sublates itself
and is rather the passing-over into nothing;
it is ceasing-to-be.
They do not sublate themselves reciprocally
[the one sublating the other externally]
but each rather sublates itself in itself
and is within it the opposite of itself.

3. Sublation of becoming

The equilibrium in which coming-to-be and ceasing-to-be are poised
is in the first place becoming itself.
But this becoming equally collects itself in quiescent unity.
Being and nothing are in it only as vanishing;
becoming itself, however, is only by virtue of their being distinguished.
Their vanishing is therefore the vanishing of becoming,
or the vanishing of the vanishing itself.
Becoming is a ceaseless unrest that collapses into a quiescent result.

This can also be expressed thus:
becoming is the vanishing of being into nothing,
and of nothing into being,
and the vanishing of being and nothing in general;
but at the same time it rests on their being distinct.
It therefore contradicts itself in itself,
because what it unites within itself is self-opposed;
but such a union destroys itself.

This result is a vanishedness, but it is not nothing;
as such, it would be only a relapse into one of
the already sublated determinations
and not the result of nothing and of being.
It is the unity of being and nothing
that has become quiescent simplicity.
But this quiescent simplicity is being,
yet no longer for itself but as determination of the whole.

Becoming, as transition into the unity of being and nothing,
a unity which is as existent
or has the shape of the one-sided
immediate unity of these moments,
is existence.

CHAPTER 2

Existence

Existence is determinate being;
its determinateness is existent determinateness, quality.
Through its quality, something is opposed to an other;
it is alterable and finite,
negatively determined not only towards an other,
but absolutely within it.
This negation in it, in contrast at first
with the finite something, is the infinite;
the abstract opposition in which these determinations appear
resolves itself into oppositionless infinity,
into being-for-itself.

The treatment of existence is therefore in three divisions:

A. existence as such
B. something and other, finitude
C. qualitative infinity.

A. EXISTENCE AS SUCH

In existence (a) as such,
its determinateness is first (b) to be distinguished as quality.
The latter, however, is to be taken in both the two determinations
of existence as reality and negation.
In these determinacies, however,
existence is equally reflected into itself,
and, as so reflected, it is posited as (c) something, an existent.

a. Existence in general

Existence proceeds from becoming. It is the simple oneness of being and
nothing. On account of this simplicity, it has the form of an immediate. Its
mediation, the becoming, lies behind it; it has sublated itself, and existence
therefore appears as a first from which the forward move is made. It is
at first in the one-sided determination of being; the other determination
which it contains, nothing, will likewise come up in it, in contrast to the
first.

It is not mere being but existence, or Dasein [in German]; according to
its [German] etymology, it is being (Sein) in a certain place (da). But the
representation of space does not belong here. As it follows upon becoming,
existence is in general being with a non-being, so that this non-being is taken
up into simple unity with being. Non-being thus taken up into being with
the result that the concrete whole is in the form of being, of immediacy,
constitutes determinateness as such.

The whole is likewise in the form or determinateness of being, since in
becoming being has likewise shown itself to be only a moment – something
sublated, negatively determined. It is such, however, for us, in our reflection;
not yet as posited in it. What is posited, however, is the determinateness as
such of existence, as is also expressed by the da (or “there”) of the Dasein. –
The two are always to be clearly distinguished. Only that which is posited in
a concept belongs in the course of the elaboration of the latter to its content.
Any determinateness not yet posited in the concept itself belongs instead to
our reflection, whether this reflection is directed to the nature of the concept
itself or is a matter of external comparison. To remark on a determinateness
of this last kind can only be for the clarification or anticipation of the whole
that will transpire in the course of the development itself. That the whole,
the unity of being and nothing, is in the one-sided determinateness of being
is an external reflection; but in negation, in something and other, and so
forth, it will become posited. – It was necessary here to call attention to
the distinction just given; but to comment on all that reflection can allow
itself, to give an account of it, would lead to a long-winded anticipation
of what must transpire in the fact itself. Although such reflections may
serve to facilitate a general overview and thus facilitate understanding,
they also bring the disadvantage of being seen as unjustified assertions,
unjustified grounds and foundations, of what is to follow. They should
be taken for no more than what they are supposed to be and should be
distinguished from what constitutes a moment in the advance of the fact
itself.

Existence corresponds to being in the preceding sphere.
But being is the indeterminate;
there are no determinations that therefore transpire in it.
But existence is determinate being, something concrete;
consequently, several determinations,
several distinct relations of its moments,
immediately emerge in it.

b. Quality

On account of the immediacy with which being and nothing are one in
existence, neither oversteps the other; to the extent that existence is existent,
to that extent it is non-being; it is determined. Being is not the universal,
determinateness not the particular. 42 Determinateness has yet to detach itself
from being; nor will it ever detach itself from it, since the now underlying
truth is the unity of non-being with being; all further determinations will
transpire on this basis. But the connection which determinateness now has
with being is one of the immediate unity of the two, so that as yet no
differentiation between the two is posited.

Determinateness thus isolated by itself, as existent determinateness, is
quality – something totally simple, immediate. Determinateness in general
is the more universal which, further determined, can be something quan-
titative as well. On account of this simplicity, there is nothing further to
say about quality as such.

Existence, however, in which nothing and being are equally contained,
is itself the measure of the one-sidedness of quality as an only immedi-
ate or existent determinateness. Quality is equally to be posited in the
determination of nothing, and the result is that the immediate or existent
determinateness is posited as distinct, reflected, and the nothing, as thus
the determinate element of determinateness, will equally be something
reflected, a negation. Quality, in the distinct value of existent, is reality;
when affected by a negating, it is negation in general, still a quality but one
that counts as a lack and is further determined as limit, restriction.

Both are an existence, but in reality, as quality with the accent on being
an existent, that it is determinateness and hence also negation is concealed;
reality only has, therefore, the value of something positive from which
negating, restriction, lack, are excluded. Negation, for its part, taken as
mere lack, would be what nothing is; but it is an existence, a quality, only
determined with a non-being.

c. Something

In existence its determinateness has been distinguished as quality; in this
quality as something existing, the distinction exists – the distinction of
reality and negation. Now though these distinctions are present in existence,
they are just as much null and sublated. Reality itself contains negation; it
is existence, not indeterminate or abstract being. Negation is for its part
equally existence, not the supposed abstract nothing but posited here as it
is in itself, as existent, as belonging to existence. Thus quality is in general
unseparated from existence, and the latter is only determinate, qualitative
being.

This sublating of the distinction is more than the mere retraction and
external re-omission of it, or a simple return to the simple beginning, to
existence as such. The distinction cannot be left out, for it is. Therefore,
what de facto is at hand is this: existence in general, distinction in it, and the
sublation of this distinction; the existence, not void of distinctions as at the
beginning, but as again self-equal through the sublation of the distinction;
the simplicity of existence mediated through this sublation. This state of
sublation of the distinction is existence’s own determinateness; existence is
thus being-in-itself; it is existent, something.

Something is the first negation of negation, as simple existent self-
reference. Existence, life, thought, and so forth, essentially take on the
determination of an existent being, a living thing, a thinking mind (“I”),
and so forth. This determination is of the highest importance if we do not
wish to halt at existence, life, thought, and so forth, as generalities – also
not at Godhood (instead of God). In common representation, something
rightly carries the connotation of a real thing. Yet it still is a very superficial
determination, just as reality and negation, existence and its determinate-
ness, though no longer the empty being and nothing, still are quite abstract
determinations. For this reason they also are the most common expressions,
and a reflection that is still philosophically unschooled uses them the most;
it casts its distinctions in them, fancying that in them it has something
really well and firmly determined. – As something, the negative of the neg-
ative is only the beginning of the subject – its in-itselfness is still quite
indeterminate. It determines itself further on, at first as existent-for-itself
and so on, until it finally obtains in the concept the intensity of the sub-
ject. At the base of all these determinations there lies the negative unity
with itself. In all this, however, care must be taken to distinguish the first
negation, negation as negation in general, from the second negation, the
negation of negation which is concrete, absolute negativity, just as the first
is on the contrary only abstract negativity.

Something is an existent as the negation of negation, for such a negation
is the restoration of the simple reference to itself – but the something is
thereby equally the mediation of itself with itself. Present in the simplicity
of something, and then with greater determinateness in being-for-itself, in
the subject, and so forth, this mediation of itself with itself is also already
present in becoming, but only as totally abstract mediation; mediation
with itself is posited in the something in so far as the latter is determined as
a simple identity. – Attention can be drawn to the presence of mediation
in general, as against the principle of the alleged bare immediacy of a
knowledge from which mediation should be excluded. But there is no
further need to draw particular attention to the moment of mediation,
since it is to be found everywhere and on all sides, in every concept.

This mediation with itself which something is in itself, when taken only
as the negation of negation, has no concrete determinations for its sides;
thus it collapses into the simple unity which is being. Something is, and
is therefore also an existent. Further, it is in itself also becoming, but a
becoming that no longer has only being and nothing for its moments.
One of these moments, being, is now existence and further an existent.
The other moment is equally an existent, but determined as the negative of
something – an other. As becoming, something is a transition, the moments
of which are themselves something, and for that reason it is an alteration – a
becoming that has already become concrete. – At first, however, something
alters only in its concept; it is not yet posited in this way, as mediated and
mediating, but at first only as maintaining itself simply in its reference to
itself; and its negative is posited as equally qualitative, as only an other in
general.

B. FINITUDE

(a) Something and other: at first they are indifferent to one another; an
other is also an immediate existent, a something; the negation thus falls
outside both. Something is in itself in contrast to its being-for-other.
But the determinateness belongs also to its in-itself, and
(b) the determination of this in-itself in turn passes over into constitution,
and this latter, as identical with determination, constitutes the imma-
nent and at the same time negated being-for-another, the limit of
something which
(c) is the immanent determination of the something itself, and the some-
thing thus is the finite.
In the first division where existence in general was considered, this exis-
tence had, as at first taken up, the determination of an existent. The
moments of its development, quality and something, are therefore of
equally affirmative determination. The present division, on the contrary,
develops the negative determination which is present in existence and was
there from the start only as negation in general. It was then the first nega-
tion but has now been determined to the point of the being-in-itself of the
something, the point of the negation of negation.
a. Something and an other
1. Something and other are, first, both existents or something.
Second, each is equally an other. It is indifferent which is named first,
and just for this reason it is named something (in Latin, when they occur
in a proposition, both are aliud, or “the one, the other,” alius alium; in the
case of an alternating relation, the analogous expression is alter alterum).
If of two beings we call the one A and the other B, the B is the one which
is first determined as other. But the A is just as much the other of the
B. Both are other in the same way. “This” serves to fix the distinction
and the something which is to be taken in the affirmative sense. But
“this” also expresses the fact that the distinction, and the privileging of one
something, is a subjective designation that falls outside the something itself.
The whole determinateness falls on the side of this external pointing; also
the expression “this” contains no distinctions; each and every something is
just as good a “this” as any other. By “this” we mean to express something
completely determinate, overlooking the fact that language, as a work of
the understanding, only expresses the universal, albeit naming it as a single
object. But an individual name is something meaningless in the sense that
it does not express a universal. It appears as something merely posited and
arbitrary for the same reason that proper names can also be arbitrarily
picked, arbitrarily given as well as arbitrarily altered. 52
Otherness thus appears as a determination alien to the existence thus
pointed at, or the other existence as outside this one existence, partly
because the one existence is determined as other only by being compared
by a Third, and partly because it is so determined only on account of the
other which is outside it, but is not an other for itself. At the same time,
as has been remarked, even for ordinary thinking every existence equally
determines itself as an other existence, so that there is no existence that
remains determined simply as an existence, none which is not outside an
existence and therefore is not itself an other.
Both are determined as something as well as other : thus they are the
same and there is as yet no distinction present in them. But this sameness of
determinations, too, falls only within external reflection, in the comparison
of the two; but the other, as posited at first, though an other with reference
to something, is other also for itself apart from the something.
Third, the other is therefore to be taken in isolation, with reference to
itself, has to be taken abstractly as the other, the t¼ ™teron of Plato who
opposes it to the one as a moment of totality, and in this way ascribes
to the other a nature of its own. Thus the other, taken solely as such, is
not the other of something, but is the other within, that is, the other of
itself. – Such an other, which is the other by its own determination, is
physical nature; nature is the other of spirit; this, its determination, is at first
a mere relativity expressing not a quality of nature itself but only a reference
external to it. But since spirit is the true something, and hence nature is
what it is within only in contrast to spirit, taken for itself the quality of
nature is just this, to be the other within, that which-exists-outside-itself (in
the determinations of space, time, matter).
The other which is such for itself is the other within it, hence the other
of itself and so the other of the other – therefore, the absolutely unequal in
itself, that which negates itself, alters itself. But it equally remains identical
with itself, for that into which it alters is the other, and this other has no
additional determination; but that which alters itself is not determined in
any other way than in this, to be an other; in going over to this other, it only
unites with itself. It is thus posited as reflected into itself with sublation of
the otherness, a self-identical something from which the otherness, which is
at the same time a moment of it, is therefore distinct, itself not appertaining
to it as something.

2. The something preserves itself in its non-being; it is essentially one
with it, and essentially not one with it. It therefore stands in reference to
an otherness without being just this otherness. The otherness is at once
contained in it and yet separated from it; it is being-for-other.
Existence as such is an immediate, bare of references; or, it is in the
determination of being. However, as including non-being within itself,
existence is determinate being, being negated within itself, and then in the
first instance an other – but, since in being negated it preserves itself at the
same time, it is only being-for-other.
It preserves itself in its non-being and is being; not, however, being
in general but being with reference to itself in contrast to its reference
to the other, as self-equality in contrast to its inequality. Such a being is
being-in-itself.

Being-for-other and being-in-itself constitute the two moments of some-
thing. There are here two pairs of determinations: (1) something and
other; (2) being-for-other and being-in-itself. The former contain the non-
connectedness of their determinateness; something and other fall apart.
But their truth is their connection; being-for-other and being-in-itself are
therefore the same determinations posited as moments of one and the same
unity, as determinations which are connections and which, in their unity,
remain in the unity of existence. Each thus itself contains within it, at the
same time, also the moment diverse from it.
Being and nothing in their unity, which is existence, are no longer
being and nothing (these they are only outside their unity); so in their
restless unity, in becoming, they are coming-to-be and ceasing-to-be. – In
the something, being is being-in-itself. Now, as self-reference, self-equality,
being is no longer immediately, but is self-reference only as the non-being
of otherness (as existence reflected into itself ). – The same goes for non-
being: as the moment of something in this unity of being and non-being:
it is not non-existence in general but is the other, and more determinedly –
according as being is at the same time distinguished from it – it is reference
to its non-existence, being-for-other.
Hence being-in-itself is, first, negative reference to non-existence; it has
otherness outside it and is opposed to it; in so far as something is in itself,
it is withdrawn from being-other and being-for-other. But, second, it has
non-being also right in it; for it is itself the non-being of being-for-other.
But being-for-other is, first, the negation of the simple reference of being
to itself which, in the first place, is supposed to be existence and something;
in so far as something is in an other or for an other, it lacks a being of
its own. But, second, it is not non-existence as pure nothing; it is non-
existence that points to being-in-itself as its being reflected into itself, just
as conversely the being-in-itself points to being-for-other.
3. Both moments are determinations of one and the same, namely of
something. Something is in-itself in so far as it has returned from the
being-for-other back to itself. But something has also a determination or
circumstance, whether in itself (here the accent is on the in) or in it; in so
far as this circumstance is in it externally, it is a being-for-other.
This leads to a further determination. Being-in-itself and being-for-other
are different at first. But that something also has in it what it is in itself and
conversely is in itself also what it is as being-for-other – this is the identity of
being-in-itself and being-for-other, in accordance with the determination
that the something is itself one and the same something of both moments,
and these are in it, therefore, undivided. – This identity already occurs
formally in the sphere of existence, but more explicitly in the treatment of
essence and later of the relations of interiority and externality, and in the
most determinate form in the treatment of the idea, as the unity of concept
and actuality. – Opinion has it that with the in-itself something lofty is
being said, as with the inner; but what something is only in itself, is also only
in it; in-itself is a merely abstract, and hence itself external determination.
The expressions: there is nothing in it, or there is something in it, imply,
though somewhat obscurely, that what is in a thing also pertains to its
in-itselfness, to its inner, true worth.
It may be observed that here we have the meaning of the thing-in-itself.
It is a very simple abstraction, though it was for a while a very important
determination, something sophisticated, as it were, just as the proposition
that we know nothing of what things are in themselves was a much valued
piece of wisdom. – Things are called “in-themselves” in so far as abstraction
is made from all being-for-other, which really means, in so far as they are
thought without all determination, as nothing. In this sense, of course, it
is impossible to know what the thing-in-itself is. For the question “what?”
calls for determinations to be produced; but since the things of which the
determinations are called for are at the same time presumed to be things-in-
themselves, which means precisely without determination, the impossibility
of an answer is thoughtlessly implanted in the question, or else a senseless
answer is given. – The thing-in-itself is the same as that absolute of which
nothing is known except that in it all is one. What there is in these things-
in-themselves is therefore very well known; they are as such nothing but
empty abstractions void of truth. What, however, the thing-in-itself in
truth is, what there basically is in it, of this the Logic is the exposition.
But in this Logic something better is understood by the in-itself than an
abstraction, namely, what something is in its concept; but this concept is
in itself concrete: as concept, in principle conceptually graspable; and, as
determined and as the connected whole of its determinations, inherently
cognizable.
Being-in-itself has at first the being-for-other as a moment standing
over against it. But positedness also comes to be positioned over against
it, and, although in this expression being-for-other is also included, the
expression still contains the determination of the bending back, which
has already occurred, of that which is not in itself into that wherein it is
positive, and this is its being-in-itself. Being-in-itself is normally to be taken
as an abstract way of expressing the concept; positing, strictly speaking, first
occurs in the sphere of essence, of objective reflection; the ground posits
that which is grounded through it; more strongly, the cause produces an
effect, an existence whose subsistence is immediately negated and which
carries the meaning that it has its substance, its being, in an other. In the
sphere of being, existence only emerges out of becoming. Or again, with
the something an other is posited; with the finite, an infinite; but the finite
does not bring forth the infinite, does not posit it. In the sphere of being,
the self-determining of the concept is at first only in itself or implicit, and
for that reason it is called a transition or passing over. And the reflecting
determinations of being, such as something and other, or finite and infinite,
although they essentially point to one another, or are as being-for-other,
also stand on their own qualitatively; the other exists; the finite, like the
infinite, is equally to be regarded as an immediate existent that stands firm
on its own; the meaning of each appears complete even without its other.
The positive and the negative, on the contrary, cause and effect, however
much they are taken in isolation, have at the same time no meaning eachE
without the other; their reflective shining in each other, the shine in each of
its other, is present right in them. – In the different cycles of determination
and especially in the progress of the exposition, or, more precisely, in the
progress of the concept in the exposition of itself, it is of capital concern
always to clearly distinguish what still is in itself or implicitly and what is
posited, how determinations are in the concept and how they are as posited
or as existing-for-other. This is a distinction that belongs only to the
dialectical development and one unknown to metaphysical philosophizing
(to which the critical also belongs); the definitions of metaphysics, like
its presuppositions, distinctions, and conclusions, are meant to assert and
produce only the existent and that, too, as existent-in-itself.
In the unity of the something with itself, being-for-other is identical with
its in-itself; the being-for-other is thus in the something. The determinate-
ness thus reflected into itself is therefore again a simple existent and hence
again a quality – determination.

b. Determination, constitution, and limit

The in-itself, in which the something is reflected into itself from its being-
for-other, no longer is an abstract in-itself but, as the negation of its
being-for-other, is mediated through this latter, which is thus its moment.
It is not only the immediate identity of the something with itself, but
the identity by virtue of which the something also has present in it what
it is in itself; the being-for-other is present in it because the in-itself is the
sublation of it, is in itself from it; but, because it is still abstract, and therefore
essentially affected with negation, it is equally affected with being-for-other.
We have here not only quality and reality, existent determinateness, but
determinateness existent-in-itself; and the development consists in positing
such determinateness as thus immanently reflected.
1. The quality which in the simple something is an in-itself essentially in
unity with the something’s other moment, its being-in-it, can be named its
determination, provided that this word is distinguished, in a more precise
signification, from determinateness in general. Determination is affirmative
determinateness; it is the in-itself by which a something abides in its
existence while involved with an other that would determine it, by which
it preserves itself in its self-equality, holding on to it in its being-for-
other. Something fulfills its determination to the extent that the further
determinateness, which variously accrues to it in the measure of its being-
in-itself as it relates to an other, becomes its filling. Determination implies
that what something is in itself is also present in it.
The determination of the human being, its vocation, is rational thought:
thinking in general is his simple determinateness; by it the human being
is distinguished from the brute; he is thinking in himself, in so far as
this thinking is distinguished also from his being-for-other, from his own
natural and sensuous being that brings him in immediate association with
the other. But thinking is also in him; the human being is himself thinking,
he exists as thinking, thought is his concrete existence and actuality; and,
further, since thinking is in his existence and his existence is in his thinking,
thinking is concrete, must be taken as having content and filling; it is rational
thought and as such the determination of the human being. But even this
determination is again only in itself, as an ought, that is to say, it is, together
with the filling embodied in its in-itself, in the form of an in-itself in general
as against the existence which is not embodied in it but still lies outside
confronting it, immediate sensibility and nature.
2. The filling of the being-in-itself with determinateness is also dis-
tinct from the determinateness which is only being-for-other and remains
outside the determination. For in the sphere of the qualitative, the dis-
tinguished terms are left, in their sublated being, also with an immediate,
qualitative being contrasting them. That which the something has in it
thus separates itself and is from this side the external existence of the some-
thing and also its existence, but not as belonging to its being-in-itself. –
Determinateness is thus constitution.
Constituted in this or that way, the something is caught up in external
influences and in external relationships. This external connection on which
the constitution depends, and the being determined through an other,
appear as something accidental. But it is the quality of the something to
be given over to this externality and to have a constitution.
In so far as something alters, the alteration falls on the side of its con-
stitution; the latter is that in the something which becomes an other. The
something itself preserves itself in the alteration; the latter affects only this
unstable surface of the something’s otherness, not its determination.
Determination and constitution are thus distinct from each other; some-
thing, according to its determination, is indifferent to its constitution.
But that which the something has in it is the middle term of this syllo-
gism connecting the two, determination and constitution. Or, rather, the
being-in-the-something showed itself to fall apart into these two extremes.
The simple middle term is determinateness as such; its identity belongs to
determination just as well as to constitution. But determination passes over
into constitution on its own, and constitution into determination. This
is implied in what has been said. The connection, upon closer consider-
ation, is this: in so far as that which something is in itself is also in it,
the something is affected with being-for-other; determination is therefore
open, as such, to the relation with other. Determinateness is at the same
time moment, but it contains at the same time the qualitative distinction
of being different from being-in-itself, of being the negative of the some-
thing, another existence. This determinateness which thus holds the other
in itself, united with the being-in-itself, introduces otherness in the latter or
in determination, and determination is thereby reduced to constitution. –
Conversely, the being-for-other, isolated as constitution and posited on its
own, is in it the same as what the other as such is, the other in it, that is,
the other of itself; but it consequently is self-referring existence, thus being-
in-itself with a determinateness, therefore determination. – Consequently,
inasmuch as the two are also to be held apart, constitution, which appears
to be grounded in something external, in an other in general, also depends
on determination, and the determining from outside is at the same time
determined by the something’s own immanent determination. And fur-
ther, constitution belongs to that which something is in itself: something
alters along with its constitution.
This altering of something is no longer the first alteration of something
merely in accordance with its being-for-other. The first was an alteration
only implicitly present, one that belonged to the inner concept; now the
alteration is also posited in the something. – The something itself is further
determined, and negation is posited as immanent to it, as its developed
being-in-itself.
The transition of determination and constitution into each other is at
first the sublation of their distinction, and existence or something in general
is thereby posited; moreover, since this something in general results from a
distinction that also includes qualitative otherness within it, there are two
somethings. But these are, with respect to each other, not just others in
general, so that this negation would still be abstract and would occur only
in the comparison of the two; rather the negation now is immanent to the
somethings. As existing, they are indifferent to each other, but this, their
affirmation, is no longer immediate: each refers itself to itself through the
intermediary of the sublation of the otherness which in determination is
reflected into the in-itselfness.
Something behaves in this way in relation to the other through itself;
since otherness is posited in it as its own moment, its in-itselfness holds
negation in itself, and it now has its affirmative existence through its
intermediary alone. But the other is also qualitatively distinguished from
this affirmative existence and is thus posited outside the something. The
negation of its other is only the quality of the something, for it is in
this sublation of its other that it is something. The other, for its part,
truly confronts an existence only with this sublation; it confronts the first
something only externally, or, since the two are in fact inherently joined
together, that is, according to their concept, their connectedness consists
in this, that existence has passed over into otherness, something into other;
that something is just as much an other as the other is. Now in so far as the
in-itselfness is the non-being of the otherness that is contained in it but
is at the same time also distinct as existent, something is itself negation,
the ceasing to be of an other in it; it is posited as behaving negatively in
relation to the other and in so doing preserving itself. This other, the in-
itselfness of the something as negation of the negation, is the something’s
being-in-itself, and this sublation is as simple negation at the same time in
it, namely, as its negation of the other something external to it. It is one
determinateness of the two somethings that, on the one hand, as nega-
tion of the negation, is identical with the in-itselfness of the somethings,
and also, on the other hand, since these negations are to each other as
other somethings, joins them together of their own accord and, since each
negation negates the other, equally separates them. This determinateness is
limit.
3. Being-for-other is indeterminate, affirmative association of something
with its other; in limit the non-being-for-other is emphasized, the quali-
tative negation of the other, which is thereby kept out of the something
that is reflected into itself. We must see the development of this concept –
a development that will rather look like confusion and contradiction.
Contradiction immediately raises its head because limit, as an internally
reflected negation of something, ideally holds in it the moments of some-
thing and other, and these, as distinct moments, are at the same time
posited in the sphere of existence as really, qualitatively, distinct.
a. Something is therefore immediate, self-referring existence and at first
it has a limit with respect to an other; limit is the non-being of the other,
not of the something itself; in limit, something marks the boundary of its
other. – But other is itself a something in general. The limit that something
has with respect to an other is, therefore, also the limit of the other as a
something; it is the limit of this something in virtue of which the something
holds the first something as its other away from itself, or is a non-being of
that something. The limit is thus not only the non-being of the other, but of
the one something just as of the other, and consequently of the something
in general.
But the limit is equally, essentially, the non-being of the other; thus,
through its limit, something at the same time is. In limiting, something
is of course thereby reduced to being limited itself; but, as the ceasing of
the other in it, its limit is at the same time itself only the being of the
something; this something is what it is by virtue of it, has its quality in it. –
This relation is the external appearance of the fact that limit is simple
negation or the first negation, whereas the other is, at the same time, the
negation of the negation, the in-itselfness of the something.
Something, as an immediate existence, is therefore the limit with respect
to another something; but it has this limit in it and is something through
the mediation of that limit, which is just as much its non-being. The limit
is the mediation in virtue of which something and other each both is and
is not.
b. Now in so far as something in its limit both is and is not, and these
moments are an immediate, qualitative distinction, the non-existence and
the existence of the something fall outside each other. Something has its
existence outside its limit (or, as representation would also have it, inside it);
in the same way the other, too, since it is something, has it outside it. The
limit is the middle point between the two at which they leave off. They have
existence beyond each other, beyond their limit; the limit, as the non-being
of each, is the other of both.
– It is in accordance with this difference of the something from its limit
that the line appears as line outside its limit, the point; the plane as plane
outside the line; the solid as solid only outside its limiting plane. – This
is the aspect of limit that first occurs to figurative representation (the self-
external-being of the concept) and is also most commonly assumed in the
context of spatial objects.
g. But further, something as it is outside the limit, as the unlimited
something, is only existence in general. As such, it is not distinguished
from its other; it is only existence and, therefore, it and its other have the
same determination; each is only something in general or each is other;
and so both are the same. But this, their at first immediate existence, is now
posited in them as limit: in it both are what they are, distinct from each
other. But it is also equally their common distinguishedness, the unity and
the distinguishedness of both, just like existence. This double identity of
the two, existence and limit, contains this: that something has existence
only in limit, and that, since limit and immediate existence are each at the
same time the negative of each other, the something, which is now only
in its limit, equally separates itself from itself, points beyond itself to its
non-being and declares it to be its being, and so it passes over into it. To
apply this to the preceding example, the one determination is this: that
something is what it is only in its limit. Therefore, the point is the limit of
line, not because the latter just ceases at the point and has existence outside
it; the line is the limit of plane, not because the plane just ceases at it; and
the same goes for the plane as the limit of solid. Rather, at the point the line
also begins; the point is its absolute beginning, and if the line is represented
as unlimited on both its two sides, or, as is said, as extended to infinity, the
point still constitutes its element, just as the line constitutes the element
of the plane, and the plane that of the solid. These limits are the principle
of that which they delimit; just as one, for instance, is as hundredth the
limit, but also the element, of the whole hundred.
The other determination is the unrest of the something in its limit in
which it is immanent, the contradiction that propels it beyond itself. Thus
the point is this dialectic of itself becoming line; the line, the dialectic of
becoming plane; the plane, of becoming total space. A second definition
is given of line, plane, and whole space which has the line come to be
through the movement of the point; the plane through the movement of
the line, and so forth. This movement of the point, the line, and so forth, is
however viewed as something accidental, or as movement only in figurative
representation. In fact, however, this view is taken back by supposing that
the determinations from which the line, and so forth, originate are their
elements and principles, and these are, at the same time, nothing else but
their limits; the coming to be is not considered as accidental or only as
represented. That the point, the line, the plane, are per se self-contradictory
beginnings which on their own repel themselves from themselves, and
consequently that the point passes over from itself into the line through
its concept, moves in itself and makes the line come to be, and so on – all
this lies in the concept of the limit which is immanent in the something.
The application itself, however, belongs to the treatment of space; as an
indication of it here, we can say that the point is the totally abstract limit,
but in a determinate existence; this existence is still taken in total abstraction,
it is the so-called absolute, that is, abstract space, the absolutely continuous
being-outside-one-another. Inasmuch as the limit is not abstract negation,
but is rather in this existence, inasmuch as it is spatial determinateness, the
point is spatial, is the contradiction of abstract negation and continuity
and is, for that reason, the transition as it occurs and has already occurred
into the line, and so forth. And so there is no point, just as there is no line
or plane.
The something, posited with its immanent limit as the contradiction of
itself by virtue of which it is directed and driven out and beyond itself, is
the finite.

c. Finitude

Existence is determinate. Something has a quality, and in this quality it is
not only determined but delimited; its quality is its limit and, affected by it,
something remains affirmative, quiescent existence. But, so developed that
the opposition of its existence and of the negation as the limit immanent to
this existence is the very in-itselfness of the something, and this is thus only
becoming in it, this negation constitutes the finitude of the something.
When we say of things that they are finite, we understand by this that
they not only have a determinateness, that their quality is not only reality
and existent determination, that they are not merely limited and as such
still have existence outside their limit, but rather that non-being consti-
tutes their nature, their being. Finite things are, but in their reference to
themselves they refer to themselves negatively – in this very self-reference
they propel themselves beyond themselves, beyond their being. They are,
but the truth of this being is (as in Latin) their finis, their end. 54 The finite
does not just alter, as the something in general does, but perishes, and its
perishing is not just a mere possibility, as if it might be without perish-
ing. Rather, the being as such of finite things is to have the germ of this
transgression 55 in their in-itselfness: the hour of their birth is the hour of
their death.

(a) The immediacy of finitude

The thought of the finitude of things brings this mournful note with it
because finitude is qualitative negation driven to the extreme, and in the
simplicity of such a determination there is no longer left to things an
affirmative being distinct from their determination as things destined to
ruin. 56 Because of this qualitative simplicity of negation that has returned
to the abstract opposition of nothing and perishing to being, finitude
is the most obstinate of the categories of the understanding; negation in
general, constitution, limit, are compatible with their other, with existence;
even the abstract nothing, by itself, is given up as an abstraction; but
finitude is negation fixed in itself and, as such, stands in stark contrast to its
affirmative. The finite thus does indeed let itself be submitted to flux; this
is precisely what it is, that it should come to an end, and this end is its only
determination. Its refusal is rather to let itself be brought affirmatively to its
affirmative, the infinite, to be associated with it; it is therefore inseparably
posited with its nothing, and thereby cut off from any reconciliation with
its other, the affirmative. The determination of finite things does not go
past their end. The understanding persists in this sorrow of finitude, for
it makes non-being the determination of things and, at the same time,
this non-being imperishable and absolute. Their transitoriness would only
pass away in their other, in the affirmative; their finitude would then be
severed from them; but this finitude is their unalterable quality, that is,
their quality which does not pass over into their other, that is, not into the
affirmative; and so finitude is eternal.
This is a very important consideration. But that the finite is absolute is
certainly not a standpoint that any philosophy or outlook, or the under-
standing, would want to endorse. The opposite is rather expressly present
in the assertion of finitude: the finite is the restricted, the perishable, the
finite is only the finite, not the imperishable; all this is immediately part
and parcel of its determination and expression. But all depends on whether
in one’s view of finitude its being is insisted on, and the transitoriness thus
persists, or whether the transitoriness and the perishing perish. The fact is
that this perishing of the perishing does not happen on precisely the view
that would make the perishing the final end of the finite. The official claim
is that the finite is incompatible with the infinite and cannot be united
with it; that the finite is absolutely opposed to the infinite. Being, absolute
being, is ascribed to the infinite. The finite remains held fast over against it
as its negative; incapable of union with the infinite, it remains absolute on
its own side; from the affirmative, from the infinite, it would receive affir-
mation and thus it would perish; but a union with the infinite is precisely
what is declared impossible. If the finite were not to persist over against
the infinite but were to perish, its perishing, as just said, would then be
the last of it – not its affirmative, which would be only a perishing of the
perishing. However, if it is not to perish into the affirmative but its end
is rather to be grasped as a nothing, then we are back at that first, abstract
nothing that itself has long since passed away.
With this nothing, however, which is supposed to be only nothing
but to which a reflective existence is nevertheless granted in thought, in
representation or in speech, the same contradiction occurs as we have
just indicated in connection with the finite, except that in the nothing
it just occurs but in the finite it is instead expressed. In the one case, the
contradiction appears as subjective; in the other, the finite is said to stand
in perpetual opposition to the infinite, in itself to be null, and to be as null
in itself. This is now to be brought to consciousness. The development
of the finite will show that, expressly as this contradiction, it collapses
internally, but that, in this collapse, it actually resolves the contradiction; it
will show that the finite is not just perishable, and that it perishes, but that
the perishing, the nothing, is rather not the last of it; that the perishing
rather perishes.

(b) Restriction and the ought

This contradiction is indeed abstractly present by the very fact that the
something is finite, or that the finite is. But something or being is no longer
posited abstractly but reflected into itself, and developed as being-in-itself
that has determination and constitution in it, or, more determinedly still,
in such a way that it has a limit within it; and this limit, as constituting
what is immanent to the something and the quality of its being-in-itself,
is finitude. It is to be seen what moments are contained in this concept of
the finite something.
Determination and constitution arose as sides for external reflection,
but determination already contained otherness as belonging to the in-itself
of something. On the one side, the externality of otherness is within the
something’s own inwardness; on the other side, it remains as otherness
distinguished from it; it is still externality as such, but in the something.
But further, since otherness is determined as limit, itself as negation of
the negation, the otherness immanent to the something is posited as the
connection of the two sides, and the unity of the something with itself (to
which both determination and constitution belong) is its reference turned
back upon itself, the reference to it of its implicitly existing determination
that in it negates its immanent limit. The self-identical in-itself thus refers
itself to itself as to its own non-being, but as negation of the negation, as
negating that which at the same time retains existence in it, for it is the
quality of its in-itselfness. Something’s own limit, thus posited by it as a
negative which is at the same time essential, is not only limit as such, but
restriction. But restriction is not alone in being posited as negative; the
negation cuts two ways, for that which it posits as negated is limit, and
limit is in general what is common to something and other, and is also
the determinateness of the in-itself of determination as such. This in-itself,
consequently, as negative reference to its limit (which is also distinguished
from it), as negative reference to itself as restriction, is the ought.
In order for the limit that is in every something to be a restriction, the
something must at the same time transcend it in itself – must refer to it
from within as to a non-existent. The existence of something lies quietly
indifferent, as it were, alongside its limit. But the something transcends
its limit only in so far as it is the sublatedness of the limit, the negative
in-itselfness over against it. And inasmuch as the limit is as restriction in
the determination itself, the something thereby transcends itself.
The ought therefore contains the double determination: once, as a deter-
mination which has an in-itselfness over against negation; and again, as a
non-being which, as restriction, is distinguished from the determination
but is at the same time itself a determination existing in itself. 57
The finite has thus determined itself as connecting determination and
limit; in this connection, the determination is the ought and the limit is the
restriction. Thus the two are both moments of the finite, and therefore both
themselves finite, the ought as well as the restriction. But only restriction is
posited as the finite; the ought is restricted only in itself, and therefore only
for us. It is restricted by virtue of its reference to the limit already immanent
within it, though this restriction in it is shrouded in in-itselfness, for
according to its determinate being, that is, according to its determinateness
in contrast to restriction, it is posited as being-in-itself.
What ought to be is, and at the same time is not. If it were, it would
not be what merely ought to be. The ought has therefore a restriction
essentially. This restriction is not anything alien; that which only ought to
be is determination now posited as it is in fact, namely as at the same time
only a determinateness.
The being-in-itself of the something is thus reduced in its determination
to the ought because the very thing that constitutes the something’s in-
itselfness is, in one and the same respect, as non-being; or again, because
in the in-itselfness, in the negation of the negation, the said being-in-itself
is as one negation (what negates) a unity with the other, and this other, as
qualitatively other, is the limit by virtue of which that unity is as reference
to it. 58 The restriction of the finite is not anything external, but the finite’s
own determination is rather also its restriction; and this restriction is both
itself and the ought; it is that which is common to both, or rather that in
which the two are identical.
But further, as “ought” the finite transcends its restriction; the same
determinateness which is its negation is also sublated, and is thus its in-
itself; its limit is also not its limit.
As ought something is thus elevated above its restriction, but conversely
it has its restriction only as ought. The two are indivisible. Something has
a restriction in so far as it has negation in its determination, and the
determination is also the being sublated of the restriction.

(c) Transition of the finite into the infinite

The ought contains restriction explicitly, for itself, and restriction contains
the ought. Their mutual connection is the finite itself, which contains them
both in its in-itself. These moments of its determination are qualitatively
opposed; restriction is determined as the negative of the ought, and the
ought equally as the negative of restriction. The finite is thus in itself the
contradiction of itself; it sublates itself, it goes away and ceases to be. 61
But this, its result, the negative as such, is (a) its very determination; for
it is the negative of the negative. So, in going away and ceasing to be,
the finite has not ceased; it has only become momentarily an other finite
which equally is, however, a going-away as a going-over into another finite,
and so forth to infinity. But, (b) if we consider this result more closely, in
its going-away and ceasing-to-be, in this negation of itself, the finite has
attained its being-in-itself; in it, it has rejoined itself. Each of its moments
contains precisely this result; the ought transcends the restriction, that is,
it transcends itself; but its beyond, or its other, is only restriction itself.
Restriction, for its part, immediately points beyond itself to its other, and
this is the ought; but this ought is the same diremption of in-itselfness and
determinateness as is restriction; it is the same thing; in going beyond itself,
restriction thus equally rejoins itself. This identity with itself, the negation
of negation, is affirmative being, is thus the other of the finite which is
supposed to have the first negation for its determinateness; this other is the
infinite.

C. INFINITY

The infinite in its simple concept can be regarded, first of all, as a fresh
definition of the absolute; as self-reference devoid of determination, it is
posited as being and becoming. The forms of existence have no place in the
series of determinations that can be regarded as definitions of the absolute,
since the forms of that sphere are immediately posited for themselves
only as determinacies, as finite in general. But the infinite is accepted
unqualifiedly as absolute, since it is explicitly determined as the negation
of the finite; the restrictedness – to which being and becoming would
somehow be susceptible even if they do not have it or exhibit it – is thereby
both explicitly referred to and denied in it.
But, in fact, by just this negation the infinite is not already free from
restrictedness and finitude. It is essential to distinguish the true concept
of infinity from bad infinity, the infinite of reason from the infinite of the
understanding. The latter is in fact a finitized infinite, and, as we shall now
discover, in wanting to maintain the infinite pure and distant from the
finite, the infinite is by that very fact only made finite.
The infinite
(a) in simple determination, is the affirmative as negation of the finite;
(b) but is thereby in alternating determination with the infinite, and is
abstract, one-sided infinite;
(c) is the self-sublation of this infinite and of the finite in one process. This
is the true infinite.
a. The infinite in general
The infinite is the negation of negation, the affirmative, being that has
reinstated itself out of restrictedness. The infinite is, in a more intense
sense than the first immediate being; it is the true being; the elevation
above restriction. At the mention of the infinite, soul and spirit light up,
for in the infinite the spirit is at home, and not only abstractly; rather, it
rises to itself, to the light of its thinking, its universality, its freedom.
What is first given with the concept of the infinite is this, that in its
being-in-itself existence is determined as finite and transcends restriction.
It is the very nature of the finite that it transcend itself, that it negate its
negation and become infinite. Consequently, the infinite does not stand
above the finite as something ready-made by itself, as if the finite stood fixed
outside or below it. Nor is it we only, as a subjective reason, who transcend
the finite into the infinite – as if, in saying that the infinite is a concept of
reason and that through reason we elevate ourselves above things temporal,
we did this without prejudice to the finite, without this elevation (which
remains external to the finite) affecting it. In so far as the finite itself is
being elevated to infinity, it is not at all an alien force that does this for it;
it is rather its nature to refer itself to itself as restriction (both restriction
as such and as ought) and to transcend this restriction, or rather, in this
self-reference, to have negated the restriction and gone above and beyond
it. It is not in the sublation of the finite in general that infinity in general
comes to be, but the finite is rather just this, that through its nature it
comes to be itself the infinite. Infinity is its affirmative determination, its
vocation, what it truly is in itself.
The finite has thus vanished into the infinite and what is, is only the
infinite.

b. Alternating determination of finite and infinite

The infinite is; in this immediacy it is at the same time the negation of
an other, of the finite. And so, as existent and at the same time as the
non-being of an other, it has fallen back into the category of the something,
of something determinate in general. More precisely: the infinite is the
existence reflected into itself which results from the mediating sublation of
determinateness in general and is consequently posited as existence distinct
from its determinateness; therefore, it has fallen back into the category of
something with a limit. In accordance with this determinateness, the finite
stands over against the infinite as real existence; they thus remain outside
each other, standing in qualitative mutual reference; the immediate being of
the infinite resurrects the being of its negation, of the finite again, which
seemed at first to have vanished into the infinite.
But the infinite and the finite are not in these referential categories only;
the two sides are further determined in addition to being as mere others
to each other. Namely, the finite is restriction posited as restriction; it is
existence posited with the determination that it passes over into what is its
in-itself and becomes infinite. Infinity is the nothing of the finite, the in-itself
that the latter ought to be, but it is this at the same time as reflected within
itself, as realized ought, as only affirmative self-referring being. In infinity
we have the satisfaction that all determinateness, alteration, all restriction
and the ought itself together with it, have vanished, are sublated, and the
nothing of the finite is posited. As this negation of the finite is the being-
in-itself determined which, as negation of negation, is in itself affirmative.
Yet this affirmation is qualitatively immediate self-reference, being; and,
because of this, the affirmative is led back to the category of being that has
the finite confronting it as an other; its negative nature is posited as existent
negation, and hence as first and immediate negation. The infinite is in this
way burdened with the opposition to the finite, and this finite, as an other,
remains a real existence even though in its being-in-itself, in the infinite,
it is at the same time posited as sublated; this infinite is that which is not
finite – a being in the determinateness of negation. Contrasted with the
finite, with the series of existent determinacies, of realities, the infinite is
indeterminate emptiness, the beyond of the finite, whose being-in-itself is
not in its existence (which is something determinate).
As thus posited over against the finite, the two connected by the quali-
tative mutual reference of others, the infinite is to be called the bad infinite,
the infinite of the understanding, for which it counts as the highest, the
absolute truth. The understanding believes that it has attained satisfaction
in the reconciliation of truth while it is in fact entangled in unreconciled,
unresolved, absolute contradictions. And it is these contradictions, into
which it falls on every side whenever it embarks on the application and
explication of these categories that belong to it, that must make it conscious
of the fact.
This contradiction is present in the very fact that the infinite remains
over against the finite, with the result that there are two determinacies.
There are two worlds, one infinite and one finite, and in their connection
the infinite is only the limit of the finite and thus only a determinate, itself
finite infinite.
This contradiction develops its content into more explicit forms. – The
finite is the real existence which persists as such even when it has gone
over into its non-being, the infinite. As we have seen, this infinite has
for its determinateness, over against the finite, only the first, immediate
negation, just as the finite, as negated, has over against this negation
only the meaning of an other and is, therefore, still a something. When,
therefore, the understanding, elevating itself above this finite world, rises
to what is the highest for it, to the infinite, the finite world remains for it
as something on this side here, and, thus posited only above the finite, the
infinite is separated from the finite and, for the same reason, the finite from
the infinite: each is placed in a different location, the finite as existence here,
and the infinite, although the being-in-itself of the finite, there as a beyond,
at a nebulous, inaccessible distance outside which there stands, enduring,
the finite.
As thus separated, they are just as much essentially connected with each
other, through the very negation that divides them. This negation connect-
ing them – these somethings reflected into themselves – is the common
limit of each over against the other; and that, too, in such a way that each
does not merely have this limit in it over against the other, but the negation
is rather the in-itselfness of each; each thus has for itself, in its separation
from the other, the limit in it. But the limit is the first negation; both are
thus limited, finite, in themselves. Yet, as each affirmatively refers itself to
itself, each is also the negation of its limit; each thus immediately repels
the negation from itself as its non-being, and, qualitatively severed from it,
posits it as an other being outside it: the finite posits its non-being as this
infinite, and the infinite likewise the finite. It is readily conceded that the
finite passes over into the infinite necessarily (that is, through its determi-
nation) and is thereby elevated to what is its in-itself, for while the finite
is indeed determined as subsistent existence, it is at the same time also a
null in itself and therefore destined to self-dissolution; whereas the infinite,
although burdened with negation and limit, is equally also determined as
the existent in-itself, so that this abstraction of self-referring affirmation is
what constitutes its determination, and hence finite existence is not present
in it. But it has been shown that the infinite itself attains affirmative being
only by the mediation of negation, as negation of negation, and that when
its affirmation thus attained is taken as just simple, qualitative being, the
negation contained in it is demoted to simple immediate negation and,
therefore, to determinateness and limit; and these, then, are excluded
from the infinite as contradicting its in-itself; they are posited as not
belonging to it but rather as opposed to its in-itself, as the finite. Since each
is in it and through its determination the positing of its other, the two are
inseparable. But this unity rests hidden in their qualitative otherness; it is
their inner unity, one that lies only at their base.
The manner of the appearance of this unity has thereby been defined.
The unity is posited in existence as a turning over or transition of the
finite into the infinite, and vice-versa; so that the infinite only emerges in
the finite, and the finite in the infinite, the other in the other; that is
to say, each arises in the other independently and immediately, and their
connection is only an external one.
The process of their transition has the following, detailed shape. We
have the finite passing over into the infinite. This passing over appears as
an external doing. In this emptiness beyond the finite, what arises? What is
there of positive in it? On account of the inseparability of the infinite and
the finite (or because this infinite, which stands apart, is itself restricted),
the limit arises. The infinite has vanished and the other, the finite, has
stepped in. But this stepping in of the finite appears as an event external
to the infinite, and the new limit as something that does not arise out of
the infinite itself but is likewise found given. And with this we are back
at the previous determination, which has been sublated in vain. This new
limit, however, is itself only something to be sublated or transcended. And
so there arises again the emptiness, the nothing, in which we find again the
said determination – and so forth to infinity.
We have before us the alternating determination of the finite and the
infinite; the finite is finite only with reference to the ought or the infinite,
and the infinite is only infinite with reference to the finite. The two are
inseparable and at the same time absolutely other with respect to each
other; each has in it the other of itself; each is thus the unity of itself and
its other, and, in its determinateness – not to be what itself and what its
other is – it is existence.
This alternating determination of self-negating and of negating the
negating is what passes as the progress to infinity, which is accepted in
so many shapes and applications as an unsurpassable ultimate at which
thought, having reached this “and so on to infinity,” has usually achieved
its end. – This progress breaks out wherever relative determinations are
pressed to the point of opposition, so that, though in inseparable unity,
each is nevertheless attributed an independent existence over against the
other. This progress is therefore the contradiction which is not resolved but
is rather always pronounced simply as present.
What we have before us is an abstract transcending which remains
incomplete because the transcending itself has not been transcended. Before
us we have the infinite; of course, this infinite is transcended, for another
limit is posited, but just because of that only a return is instead made back
to the finite. This bad infinite is in itself the same as the perpetual ought;
it is indeed the negation of the finite, but in truth it is unable to free
itself from it; the finite constantly resurfaces in it as its other, since this
infinite only is with reference to the finite, which is its other. The progress
to infinity is therefore only repetitious monotony, the one and the same
tedious alternation of this finite and infinite.
The infinity of the infinite progress remains burdened by the finite as
such, is thereby restricted, and is itself finite. In fact, however, it is thereby
posited as the unity of the finite and the infinite. Only, this unity is not
reflected upon. Yet it alone rouses the finite in the infinite, and the infinite
in the finite; it is, so to speak, the impulse driving the infinite progress. This
progress is the outside of this unity at which representation remains fixated –
fixated at that perennial repetition of one and the same alternation; at the
empty unrest of a progression across the limit towards the infinite which, in
this infinity, finds a new limit but is just as unable to halt at it as it is at the
infinite. This infinite has the rigid determination of a beyond that cannot
be attained, for the very reason that it ought not be attained, since the
determinateness of the beyond, of an existent negation, has not been let go.
In this determination, the infinite has the finite as a this-side over against it –
a finite that is likewise unable to raise itself up to the infinite just because
it has this determination of an other, that is, of an existence that perennially
regenerates itself in that beyond precisely by being different from it.
In this reciprocal determination of the finite and the infinite alternating
back and forth as just indicated, the truth of these two is already implicitly
present in itself, and all that is needed is to take up what is there. This back
and forth movement constitutes the external realization of the concept in
which the content of the latter is posited, but externally, as a falling out of
the two; all that is needed is the comparing of these two different moments
in which the unity is given which the concept itself gives. “Unity of the
finite and the infinite” – as has often been already noted but must especially
be kept in mind at this juncture – is the uneven expression for the unity as
it is in truth; but also the removal of this uneven determination must be
found in the externalization of the concept that lies ahead of us.
Taken in their first, only immediate determination, the infinite is the
transcending of the finite; according to its determination, it is the negation
of the finite; the finite, for its part, is only that which must be transcended,
the negation in it of itself, and this is the infinite. In each, therefore, there is
the determinateness of the other, whereas, according to the viewpoint of the
infinite progression, the two should be mutually excluded and would have
to follow one another only alternately; neither can be posited and grasped
without the other, the infinite without the finite, the finite without the
infinite. In saying what the infinite is, namely the negation of the finite,
the finite itself is said also; it cannot be avoided in the determination
of the infinite. One need only know what is being said in order to find
the determination of the finite in the infinite. Regarding the finite, it
is readily conceded that it is the null; this very nothingness is however
the infinite from which it is inseparable. – Understood in this way, they
may seem to be taken according to the way each refers to its other. Taken
without this connecting reference, and thus joined only through an “and,”
they subsist independently, each only an existent over against the other.
We have to examine how they would be constituted in this way. The
infinite, thus positioned, is one of the two; but, as only one of them, it
is itself finite, it is not the whole but only one side; it has its limit in
that which stands over against it; and so it is the finite infinite. We have
before us only two finites. The finitude of the infinite, and therefore its
unity with the finite, lies in the very fact that it is separated from the finite
and placed, consequently, on one side. – The finite, for its part, removed
from the infinite and positioned for itself, is this self-reference in which the
relativity, its dependence and transitoriness, are removed; it is the same
self-subsistence and self-affirmation which the infinite is presumed to be.
The two pathways of consideration, even though they seem at first to
have each a different determinateness for their point of departure – the
former inasmuch as it assumes it to be only the reference of infinite and
finite to each other, of each to the other; and the latter their complete
separation from each other – yield one and the same result. The infinite
and the finite, taken together as referring to each other in a connection
which is presumed external but is in fact essential to them (for without it,
neither is what it is), each contains its other in its own determination, just
as, when each is taken for itself, when looked at on its terms, each has the
other present in it as its own moment.
This yields, then, the scandalous unity of the finite and the infinite –
the unity which is itself the infinite that embraces both itself and the
finite – the infinite, therefore, understood in a sense other than when the
finite is separated from it and placed on the other side from it. Since they
must now also be distinguished, each is within it, as just shown, itself
the unity of both; there are thus two such unities. The common element,
the unity of both determinacies, as such a unity, posits them at first as
negated, for each is to be what it is in being distinguished; in their unity,
therefore, they lose their qualitative nature – an important reflection for
countering the incorrigible habit of representing the infinite and the finite,
in their unity, as still holding on to the quality that they would have when
taken apart from each other; of seeing in that unity, therefore, nothing
except contradiction, and not also the resolution of the contradiction by
the negation of the qualitative determinateness of each. And so is the unity
of the infinite and the finite, at first simple and universal, falsified.
But further, since the two are now to be taken also as distinguished, the
unity of the infinite which is itself both of these moments is determined
differently in each. The infinite, determined as such, has in it the finitude
which is distinct from it; in this unity, the infinite is the in-itself while the
finite is only determinateness, the limit in the infinite. But such a limit is
the absolute other of the infinite, its opposite. The infinite’s determination,
which is the in-itself as such, is corrupted by being saddled with a quality
of this sort; the infinite is thus a finitized infinite. Likewise, since the finite
is as such only the non-in-itself but equally has its opposite in it by virtue
of the said unity, it is elevated above its worth and, so to speak, infinitely
elevated; it is posited as the infinitized finite.
Likewise, just as the simple unity of infinite and finite was falsified before
by the understanding, so too is the double unity. Here also this happens
because the infinite is taken in one of the two unities not as negated but,
rather, as the in-itself in which, therefore, determinateness and restriction
should not be posited, for they would debase and corrupt it. Conversely,
the finite is equally held fixed as not negated, although null in itself; so
that, in combination with the infinite, it is elevated to what it is not and is
thereby infinitized notwithstanding its determination that has not vanished
but is rather perpetuated.
The falsification that the understanding perpetrates with respect to the
finite and the infinite, of holding their reciprocal reference fixed as quali-
tative differentiation, of maintaining that their determination is separate,
indeed, absolutely separate, comes from forgetting what for the under-
standing itself is the concept of these moments. According to this concept,
the unity of the finite and the infinite is not an external bringing together of
them, nor an incongruous combination that goes against their nature, one
in which inherently separate and opposed terms that exist independently
and are consequently incompatible, would be knotted together. Rather,
each is itself this unity, and this only as a sublating of itself in which neither
would have an advantage over the other in in-itselfness and affirmative
existence. As has earlier been shown, finitude is only as a transcending
of itself; it is therefore within it that the infinite, the other of itself, is
contained. Similarly, the infinite is only as the transcending of the finite;
it therefore contains its other essentially, and it is thus within it that it is
the other of itself. The finite is not sublated by the infinite as by a power
present outside it; its infinity consists rather in sublating itself.
This sublating is not, consequently, alteration or otherness in general,
not the sublating of something. That into which the finite is sublated is
the infinite as the negating of finitude. But the latter has long since been
only existence, determined as a non-being. It is only the negation, therefore,
that in the negation sublates itself. Thus infinity is determined on its side
as the negative of the finite and thereby of determinateness in general, as
an empty beyond; its sublating of itself into the finite is a return from an
empty flight, the negation of the beyond which is inherently a negative.
Present in both, therefore, is the same negation of negation. But this
negation of negation is in itself self-reference, affirmation but as turning
back to itself, that is, through the mediation that the negation of negation is.
These are the determinations that it is essential to bring to view; the second
point, however, is that in the infinite progression they are also posited, and
how they are posited therein, namely, not in their ultimate truth.
First, both are negated in that progression, the infinite as well as the finite;
both are equally transcended. Second, they are also posited as distinct, one
after the other, each positive for itself. We sort out these two determinations
while comparing them, just as in the comparison (in an external comparing)
we have separated the two ways of considering them: the finite and the
infinite as referring to one another, and each taken for itself. The infinite
progression, however, says more than this. Also posited in it, though at
first still only as transition and alternation, is the connectedness of the terms
being distinguished. We now only need to see, in one simple reflection,
what is in fact present in it.
In the first place, the negation of the finite and the infinite which is
posited in the infinite progression can be taken as simple, and hence as
mutual externality, only a following of one upon the other. Starting from
the finite, the limit is thus transcended, the finite negated. We now have
its beyond, the infinite, but in this the limit rises up again; so we have the
transcending of the infinite. This twofold sublation is nonetheless partly
only an external event and an alternating of moments in general, and partly
still not posited as one unity; each of these moves beyond is an independent
starting point, a fresh act, so that the two fall apart. – But, in addition, their
connection is also present in the infinite progression. The finite comes first;
then there is the transcending of it, and this negative, or this beyond of the
finite, is the infinite; third, this negation is transcended in turn, a new limit
comes up, a finite again. – This is the complete, self-closing movement that
has arrived at that which made the beginning; what emerges is the same as
that from which the departure was made, that is, the finite is restored; the
latter has therefore rejoined itself, in its beyond has only found itself again.
The same is the case regarding the infinite. In the infinite, in the beyond
of the limit, only a new limit arises which has the same fate, namely, that
as finite it must be negated. Thus what is again at hand is the same infinite
that just now disappeared in the new limit; by being sublated, by traversing
the new limit, the infinite has not therefore advanced one jot further: it
has distanced itself neither from the finite (for the finite is just this, to pass
over into the infinite), nor from itself, for it has arrived at itself.
Thus the finite and the infinite are both this movement of each returning
to itself through its negation; they are only as implicit mediation, and the
affirmative of each contains the negative of each, and is the negation of the
negation. – They are thus a result and, as such, not in the determination
that they had at the beginning: neither is the finite an existence on its side
nor the infinite an existence or a being-in-itself beyond that existence, that
is, beyond existence in the determination of finitude. The understanding
strongly resists the unity of the finite and the infinite only because it pre-
supposes restriction and finitude to remain, like being-in-itself, constants.
It thereby overlooks the negation of both which is in fact present in the
infinite progression, just as it equally overlooks that the two occur in this
progression only as moments of a whole – that each emerges only through
the mediation of its opposite but, essentially, equally by means of the
sublation of its opposite.
If this immanent turning back has for the moment been reckoned to
be just as much the turning back of the finite to itself and of the infinite
to itself, noticeable in this very result is an error connected with the one-
sidedness just criticized: the finite and then the infinite is each taken as
the starting point, and only in this way two results ensue. But it is a matter
of total indifference which is taken as the starting point and, with this,
the distinction caused by the duality of results dissolves of itself. This is
likewise posited in the line of the infinite progression, open-ended on both
sides, wherein each of the moments recurs in equal alternation, and it is
totally extraneous at which position the progression is arrested and taken
as beginning. – The moments are distinguished in the progression but each
is equally only moment of the other. Since both, the finite and the infinite,
are themselves moments of the progress, they are jointly the finite, and,
since they are equally jointly negated in it and in the result, this result as
the negation of their joint finitude is called with truth the infinite. Their
distinction is thus the double meaning which they both have. The finite has
the double meaning, first, of being the finite over against the infinite which
stands over against it, and, second, of being at the same time the finite
and the infinite over against the infinite. Also the infinite has the double
meaning of being one of the two moments (it is then the bad infinite)
and of being the infinite in which the two moments, itself and its other,
are only moments. Therefore, as in fact we now have it, the nature of the
infinite is that it is the process in which it lowers itself to be only one of its
determinations over against the finite and therefore itself only one of the
finites, and elevates this distinction of itself and itself to be self-affirmation
and, through this mediation, the true infinite.
This determination of the true infinite cannot be captured in the already
criticized formula of a unity of the finite and the infinite; unity is abstract,
motionless self-sameness, and the moments are likewise unmoved beings.
But, like both its moments, the infinite is rather essentially only as becoming,
though a becoming now further determined in its moments. Becoming has
for its determinations, first, abstract being and nothing; as alteration, it has
existence, something and other; now as infinite, it has finite and infinite,
these two themselves as in becoming.
This infinite, as being-turned-back-unto-itself, as reference of itself
to itself, is being – but not indeterminate, abstract being, for it is
posited as negating the negation; consequently, it is also existence or
“thereness,” 62 for it contains negation in general and consequently deter-
minateness. It is, and is there, present, before us. Only the bad infinite is
the beyond, since it is only the negation of the finite posited as real and,
as such, it is abstract first negation; thus determined only as negative, it
does not have the affirmation of existence in it; held fast only as something
negative, it ought not to be there, it ought to be unattainable. However, to
be thus unattainable is not its grandeur but rather its defect, which is at
bottom the result of holding fast to the finite as such, as existent. It is the
untrue which is the unattainable, and what must be recognized is that such
an infinite is the untrue. – The image of the progression in infinity is the
straight line; the infinite is only at the two limits of this line, and always
only is where the latter (which is existence) is not but transcends itself, in
its non-existence, that is, in the indeterminate. As true infinite, bent back
upon itself, its image becomes the circle, the line that has reached itself,
closed and wholly present, without beginning and end.
True infinity, thus taken in general as existence posited as affirmative in
contrast to abstract negation, is reality in a higher sense than it was earlier
as simply determined; it has now obtained a concrete content. It is not
the finite which is the real, but rather the infinite. Thus reality is further
determined as essence, concept, idea, and so forth. In connection with the
more concrete, it is however superfluous to repeat such earlier and more
abstract categories as reality, and to use them for determinations more
concrete than they are by themselves. Such a repetition, as when it is said
that essence, or that the concept, is real, has its origin in the fact that to
uneducated thought the most abstract categories such as being, existence,
reality, finitude, are the most familiar.
The more immediate occasion, however, for recalling here the categories
of reality is that the negation, against which reality is the affirmative, is here
the negation of negation, and consequently itself posited over against that
reality which finite existence is. – Negation is thus determined as ideality;
the idealized h is the finite as it is in the true infinite – as a determination, a
content, a distinct but not a subsistent existent, a moment rather. Ideality has
this more concrete signification which is not fully expressed through the
negation of finite existence. – As regards reality and ideality, the opposition
The ideal [das Ideale] has a broader meaning (such as of the beautiful and its associations) than
the idealized [das Ideelle]. The former does not belong here yet, and for this reason the expression
idealized is being used. There is no such distinction made in language usage for “reality”; in German
reelle and reale are used as roughly synonymous and no interest is served in nuancing the two in
some sort of opposition.
of finite and infinite is, however, so grasped that the finite assumes the value
of “the real,” whereas the infinite that of “the idealized”; in the same way,
further on, also the concept is regarded as an idealization, that is, as a mere
idealization, in contrast to existence in general, which is regarded as “the
real.” When contrasted in this way, it is of course of no use to have reserved
for the said concrete determination of negation the distinctive expression
of “idealization”; in that opposition of finite and infinite, we are back to
the one-sidedness of the abstract negative characteristic of the bad infinite
and still fixed in the affirmative existence of the finite.

Transition

Ideality can be called the quality of the infinite;
but it is essentially the process of becoming,
and hence a transition, like the transition
of becoming into existence.
We must now explicate this transition.
This immanent turning back, as the sublating of finitude,
of finitude as such and equally of the negative finitude that
only stands opposite to it, is only negative finitude,
is self-reference, being.
Since there is negation in this being, the latter is existence;
but, further, since the negation is essentially
negation of the negation, self-referring negation,
it is the existence that carries the name of being-for-itself.

CHAPTER 3

Being-for-itself

In being-for-itself,
qualitative being is brought to completion;
it is infinite being;
the being of the beginning is void of determination;
existence is sublated but only immediately sublated being;
it thus contains, to begin with,
only the first negation, itself immediate;
being is of course retained as well,
and the two are united in existence in simple unity;
for this reason, however, each is in itself still unlike the other,
and their unity is still not posited.
Existence is therefore the sphere of differentiation,
of dualism, the domain of finitude.
Determinateness is determinateness as such;
being which is relatively, not absolutely, determined.
In being-for-itself, the distinction
between being and determinateness,
or negation, is posited and equalized.
Quality, otherness, limit, as well as reality, in-itselfness,
ought, and so forth, are the incomplete configurations of negation in being
which are still based on the differentiation of the two.
But since in finitude negation has passed over into infinity,
in the posited negation of negation,
negation is simple self-reference
and in it, therefore, the equalization with being:
absolutely determinate being.

First, being-for-itself is immediately
an existent-for-itself, the one.

Second, the one passes over into a multiplicity of ones,
repulsion or the otherness of the one
which sublates itself into its ideality, attraction.

Third, we have the alternating determination of repulsion and attraction
in which the two sink into a state of equilibrium;
and quality, driven to a head in being-for-itself,
passes over into quantity.

A. BEING-FOR-ITSELF AS SUCH

The general concept of being-for-itself has come to light.
The justification for using the expression “being-for-itself”
for that concept would depend on showing that the representation
associated with the expression corresponds to the concept.
So indeed it appears to do. We say that something is for itself
inasmuch as it sublates otherness,
sublates its connection and community with other,
has rejected them by abstracting from them.
The other is in it only as something sublated, as its moment;
being-for-itself consists in
having thus transcended limitation, its otherness;
it consists in being, as
this negation, the infinite turning back into itself.
In representing to itself an intended object
which it feels, or intuits, and so forth,
consciousness already contains in itself as consciousness
the determination of being-for-itself;
that is, it has in it the content of that object,
which is thus an idealization;
even as it intuits, or in general becomes involved in
the negative of itself, in the other, it abides with itself.
Being-for-itself is the polemical,
negative relating to the limiting other
and, through this negation of the other,
is being-reflected-within-itself;
even though, side by side with this
immanent turning back of consciousness
and the ideality of its object,
the reality of this object is also retained,
for the object is at the same time
known as an external existence.
Consciousness is thus phenomenal,
or it is this dualism:
on the one side, it knows an external object
which is other than it;
on the other side, it is for-itself,
has this intended object in it as idealized,
abides not only by this other
but therein abides also with itself.
Self-consciousness, on the contrary, is
being-for-itself brought to completion and posited;
the side of reference to another, to an external object, is removed.
Self-consciousness is thus the nearest example of the presence of infinity;
granted, of a still abstract infinity,
but one which is of a totally different, concrete determination
than the being-for-itself in general,
whose infinity still has only qualitative determinateness.

a. Existence and being-for-itself

As already mentioned, being-for-itself is
infinity that has sunk into simple being;
it is existence in so far as in the now posited form of the immediacy
of being the negative nature of infinity,
which is the negation of negation,
is only as negation in general,
as infinite qualitative determinateness.
But in such a determinateness, wherein it is existence,
being is at once also distinguished
from this very being-for-itself
which is such only as infinite
qualitative determinateness;
nevertheless, existence is at the same time
a moment of being-for-itself,
for the latter certainly contains
being affected by negation.
So the determinateness which in existence as such is
an other, and a being-for-other,
is bent back into the infinite unity of being-for-itself,
and the moment of existence is present
in the being-for-itself as being-for-one.

b. Being-for-one

This moment gives expression to how the finite is
in its unity with the infinite or as an idealization.
Being-for-itself does not have negation in it as
a determinateness or limit,
and consequently also not as reference
to an existence other than it.
Although this moment is now being
designated as being-for-one,
there is yet nothing at hand for which it would be;
there is not the one of which it would be the moment.
There is in fact nothing of the sort
yet fixed in being-for-itself;
that for which something (and there is no something here)
would be, what the other side in general should be,
is likewise a moment, itself only being-for-one, not yet a one.
What we have before us, therefore,
is still an undistinguishedness of two sides
that may suggest themselves in the being-for-one;
there is only one being-for-another,
and since this is only one being-for-another,
it is also only being-for-one;
there is only the one ideality,
of that for which or in which
there should be a determination as moment,
and of that which should be the moment in it.
Being-for-one and being-for-itself do not therefore
constitute two genuine determinacies,
each as against the other.
Inasmuch as the distinction is momentarily assumed
and we speak of a-being-for-itself,
it is this very being-for-itself,
as the sublated being of otherness,
that refers itself to itself as to the sublated other,
is therefore for-one;
in its other it refers itself only to itself.
An idealization is necessarily for-one,
but it is not for an other;
the one, for which it is, is only itself.
The “I,” therefore, spirit in general,
or God, are idealizations,
because they are infinite;
as existents which are for-themselves, however,
they are not ideationally different from that which is for-one.
For if they were different, they would be only immediate,
or, more precisely, they would only be
existence and a being-for-another;
for if the moment of being for-one did
not attach to them,
it is not they themselves
but an other that would be
that which is for them.
God is therefore for himself,
in so far he is himself
that which is for him.

Being-for-itself and being-for-one are not, therefore,
diverse significations of ideality
but essential, inseparable, moments of it.

c. The one

Being-for-itself is the simple unity of
itself and its moments, of the being-for-one.
There is only one determination present,
the self-reference itself of the sublating.
The moments of being-for-itself have sunk into
an indifferentiation which is immediacy or being,
but an immediacy that is based on
the negating posited as its determination.
Being-for-itself is thus an existent-for-itself,
and, since in this immediacy its inner meaning vanishes,
it is the totally abstract limit of itself: the one.

Attention may be drawn in advance to the difficulties that lie ahead
in the exposition of the development of the one,
and to the source of these difficulties.
The moments that constitute the concept of
the one as being-for-itself
occur in it one outside the other;
they are
(1) negation in general;
(2) two negations that are, therefore,
(3) the same,
(4) absolutely opposed;
(5) self-reference, identity as such;
(6) negative reference which is nonetheless self-reference.
These moments occur here apart because
the form of immediacy, of being, enters into
the being-for-itself as existent-for-itself;
because of this immediacy, each moment is posited
as a determination existent on its own,
and yet they are just as inseparable.
Hence, of each determination the opposite must equally be said;
it is this contradiction that causes the difficulty
that goes with the abstract nature of the moments.

B. THE ONE AND THE MANY

The one is the simple reference of being-for-itself to itself in which its
moments have fallen together – in which, therefore, being-for-itself has the
form of immediacy and its moments, therefore, are now there as existents.
As the self-reference of the negative, the one is a determining – and, as
self-reference, it is infinite self-determining. However, because of the present
immediacy, these distinctions are no longer only moments of one and the
same self-determination but are at the same time posited as existents. The
ideality of the being-for-itself as a totality thus turns at first into reality – a
reality, moreover, of the most fixed and abstract kind, as a one. In the one,
the being-for-itself is the posited unity of being and existence, as the absolute
union of the reference to another and the reference to itself; but also the
determinateness of being then enters into opposition to the determination
of the infinite negation, to self-determination, so that what the one is in
itself, it is that now only in it, and the negative consequently is an other
distinct from it. What shows itself to be present as distinct from the one is
the one’s own self-determining; its unity with itself, as thus distinct from
itself, is demoted to reference, and, as negative unity, it is negation of itself
as other, the excluding of the one as an other from itself, from the one.

a. The one within

Within it, the one just is; this, its being, is not an existence, not a deter-
mination as reference to an other, not a constitution; it is rather its having
negated this circle of categories. The one is not capable, therefore, of
becoming any other; it is unalterable.
It is indeterminate, yet no longer like being; its indeterminateness is the
determinateness of self-reference, absolutely determined being; posited in-
itselfness. As negation which, in accordance with its concept, is self-referring,
it has distinction in it: it directs away from itself towards another, but this
direction is immediately reversed, because, according to this moment of
self-determining, there is no other to which it would be addressed, and the
directing reverts back to itself.
In this simple immediacy, even the mediation of existence and ideality,
and with it all diversity and manifoldness, have vanished. In the one there is
nothing; this nothing, the abstraction of self-reference, is here distinguished
from the in-itselfness of the one; it is a posited nothing, for this in-itselfness
no longer has the simplicity of the something, but, as mediation, has
rather the determination of being concrete; taken in abstraction, it is
indeed identical with the one, but different from its determination. So this
nothing, posited as in the one, is the nothing as the void. – The void is thus
the quality of the one in its immediacy.

b. The one and the void

The one is the void as the abstract self-reference of negation. But the void,
as nothing, is absolutely diverse from the simple immediacy of the one,
from the being of the latter which is also affirmative, and because the two
stand in one single reference, namely to the one, their diversity is posited;
however, as distinct from the affirmative being, the nothing stands as void
outside the one as existent.
Being-for-itself, determined in this way as the one and the void, has
again acquired an existence. – The one and the void have their negative
self-reference as their common and simple terrain. The moments of being-
for-itself come out of this unity, become external to themselves; for through
the simple unity of the moments the determination of being comes into
play, and the unity itself thus withdraws to one side, is therefore lowered
to existence, and there it is confronted by its other determination standing
over against it, negation as such and likewise as the existence of the nothing,
as the void.

c. Many ones

Repulsion

The one and the void constitute the first existence of being-for-itself. Each
of these moments has negation for its determination, and is posited at the
same time as an existence. In accordance with this determination, the one
and the void are each the reference of negation to negation as of an other
to its other: the one is negation in the determination of being; the void,
negation in the determination of non-being. Essentially, however, the one
is only self-reference as referring negation, that is, it is itself the same as the
void outside it is supposed to be. Both are, however, also posited as each
an affirmative existence – the one as being-for-itself as such, the other as
indeterminate existence in general – and each as referring to the other as
to an other existence. Essentially, however, the being-for-itself of the one
is the ideality of the existence and of the other; it does not refer to an
other but only to itself. But inasmuch as the being-for-itself is fixed as the
one, as existent for itself, as immediately present, its negative reference to
itself is at the same time reference to an existent; and since the reference
is just as much negative, that to which the being-for-itself refers remains
determined as an existence and as an other; as essentially self-reference, the
other is not indeterminate negation like the void, but is likewise a one. The
one is consequently a becoming of many ones.

Strictly speaking, however, this is not just a becoming; for becoming is
a transition of being into nothing; the one, by contrast, becomes only a
one. The one, as referred to, contains the negative as reference; it has this
reference, therefore, in it. Hence, instead of a becoming, the one’s own
immanent reference is, first, present; and, second, since this reference is
negative and the one is at the same time an existent, the one repels itself
from itself. This negative reference of the one to itself is repulsion.
This repulsion, as thus the positing of many ones but through the one
itself, is the one’s own coming-forth-from-itself, but to such outside it as
are themselves only ones. This is repulsion according to the concept, as it
exists implicitly in itself. The second repulsion is distinguished from it. It
is the one that first occurs to the representation of external reflection, not
as the generation of ones but only as the mutual holding off of ones which
are presupposed as already there. To be seen now is how the first repulsion
that exists in itself determines itself as the second, the external repulsion.
We must first establish the determinations that the many ones have as
such. The becoming of the many, or their being produced, immediately
vanishes as the product of a positing; what is produced are the ones, not
for another, but as infinitely referring to themselves. The one repels only
itself from itself; it does not come to be but it already is; that which is
represented as the repelled is equally a one, an existent; repelling and being
repelled applies in like manner to both, and makes no difference.
The ones are thus presupposed with respect to each other – posited through
the repulsion of the one from itself; pre-supposed, posited as non-posited;
their being-posited is sublated, they are existents with respect to each other,
such as refer only to themselves.
Thus plurality appears not as an otherness, but as a determination com-
pletely external to the one. The one, in repelling itself, remains reference to
itself, just like that which is taken as repelled at the start. That the ones are
other to one another, that they are brought together in the determinateness
of plurality, does not therefore concern the one. If the plurality were a
reference of the ones to one another, the ones would then limit each other
and would have the being-for-other affirmatively in them. Their connect-
ing reference (and this they have through their unity which is in itself ),
as posited here, is determined as none; it is again the previously posited
void. This void is their limit, but an external limit in which they are not
supposed to be for one another. The limit is that in which the limited are
just as much as are not; but the void is determined as pure non-being, and
this alone constitutes the limit of the ones.
The repulsion of the one from itself is the making explicit of what the
one is implicitly in itself; but, thus laid out as one-outside-the-other, infinity
is here an infinity that has externalized itself, and this it has done through
the immediacy of the infinite, of the one. Infinity is just as much the simple
reference of the one to the one as, on the contrary, the one’s absolute lack
of reference; it is the former according to the simple affirmative reference
of the one to itself; it is the latter according to the same reference as
negative. Or again, the plurality of the ones is the one’s own positing of
the one; the one is nothing but the negative reference of the one to itself,
and this reference – hence the one itself – is the plural one. But equally,
plurality is utterly external to the one, for the one is precisely the sublating
of otherness; repulsion is its self-reference and simple equality with itself.
The plurality of the ones is infinity as a contradiction that unconstrainedly
produces itself.

C. REPULSION AND ATTRACTION

a. Exclusion of the one

The many ones are each a being; their existence or their reference to
one another is a non-reference, it is external to them: the abstract void.
But they themselves are now this negative reference to themselves as to
existent others: the demonstrated contradiction, the infinity posited in the
immediacy of being. With this, repulsion now finds immediately before it
that which is repelled by it. In this determination, it is an excluding; the one
repels from itself only the many not generated by it, the ones not posited
by it. This repelling is mutual or from all sides – relative, limited by the
being of the ones.
Plurality is not at first posited otherness; limit is only the void, only that
in which the ones are not. But in the limit they also are; they are in the
void, or their repulsion is their common connecting reference.
This mutual repulsion is the posited existence of the many ones; it is not
their being-for-itself, in accordance with which they would be distinguished
as many only in a third, but is rather their own distinguishing which
preserves them. – They mutually negate themselves, posit one another as
being only for-one. But at the same time they negate this being only for-
one just as much; they repel the ideality that they have and are. – So the
moments which in ideality are absolutely united come apart. In its being-
for-itself, the one is also for-one; but this one, for which it is, is itself; its
distinguishing from itself is immediately sublated. But in the plurality the
distinguished one has a being; the being-for-one as has been determined
in exclusion is therefore a being-for-other. Each thus comes to be repelled
by an other, is sublated and made into a one which is not for itself but
for-one, and an other one at that.
The being-for-itself of the many ones thus shows itself to be their self-
preservation through the mediation of their mutual repulsion in which
they sublated themselves reciprocally and posit the others as mere being-
for-another. But the self-preservation consists at the same time in repelling
this ideality and positing the ones as not being for-an-other. This self-
preservation of the ones through their negative reference to one another is,
however, rather their dissolution.
The ones not only are but maintain themselves through their recipro-
cal exclusion. First, it is in their being, and indeed their being-in-itself as
contrasted with their reference to the other, that they should now have a
firm point of support for their diversity as against their being negated; this
in-itselfness rests on their being ones. But they all are this; in their being-in-
itself, instead of having there their firm point of support for their diversity,
they are all the same. Second, their existence and their way of relating
to one another, that is, their positing themselves as one, is their recipro-
cal negating; this, however, is likewise one and the same determination of
all through which they therefore posit themselves as identical; just as, by
being in themselves the same, the ideality that should be posited in them
through others is their own, and they thus repel just as little. – According
to their being and positing, they are, consequently, only one affirmative
unity.
This consideration regarding the ones – that from either side of their
determination, whether they just are or refer to one another, they show
themselves to be only one and the same, indistinguishable – is a comparison
that belongs to us. – Also to be seen, therefore, is what is posited in them
in their mutual reference itself. – They are – this much is presupposed
in this reference – and they are only inasmuch as they negate themselves
reciprocally and at the same time keep away this ideality, their being
negated, from themselves, that is, they negate the reciprocal negating.
But they are only inasmuch as they negate, and so, since their reciprocal
negating is negated, their being is negated. To be sure, since they are,
nothing would be negated through this negating which for them is only
something external; this negating of the other rebounds off them, coming
their way only by striking their surface. And yet, they turn back upon
themselves only by negating the others; they are only as this mediation,
this turning back of theirs is their self-preservation and their being-for-
itself. Since their negating is ineffectual because of the resistance offered by
the others, whether as existents or as negating, they do not return back to
themselves, do not preserve themselves, and so are not.
It was previously remarked that the ones themselves are each a one like
any other. 77 This is not just a matter of our connecting them by way of
reference, of bringing them together externally; repulsion is itself a referring;
the one that excludes the ones refers itself to them, to the ones, that is,
to itself. The negative relating of the ones to one another is consequently
only a coming-together-with-oneself. This identity in which their repelling
crosses over is the sublation of their diversity and externality which they
should have rather asserted with respect to each other by excluding each
other.
This self-positing-in-a-one of the many ones is attraction.

b. The one one of attraction

Repulsion is the fragmentation of the one, first into the many of which it
is the negative relating, since they presuppose each other as each existent;
it is only the ought of ideality; this ideality will, however, be realized in
attraction. Repulsion passes over into attraction, the many ones into one
one. Both, repulsion and attraction, are at first distinguished from each
other, repulsion as the reality of the ones, attraction as their posited ideality.
Attraction refers to repulsion by having it for a presupposition. Repulsion
delivers the material for attraction. If there were no ones, there would
be nothing to attract; the representation of continuing attraction, of the
consumption of the ones, presupposes an equally continuing generation
of the ones; the sense representation of spatial attraction gives continuity
to the flow of ones to be attracted; to replace the atoms that vanish at the
point of attraction, another multitude comes forth from the void, infinitely
if one so wishes. If attraction were represented as accomplished, that is,
the many as brought to the point of the one one, the result would be just
an inert one, no longer any attraction. The ideality immediately present in
attraction still also has in it the determination of the negation of itself, the
many ones to which it refers; attraction is inseparable from repulsion.
To attract pertains at first in equal measure to each of the many ones as
immediately present; none has advantage over an other; what would result
then is an equilibrium in the attraction, or more precisely, an equilibrium
in the attraction and the repulsion themselves, and an inert state of rest
without any ideality present there. But there can be no question here of
any such immediately present one taking precedence over another, for
this would presuppose a determinate distinction between them; attraction
is rather the positing of the given lack of distinction among the ones.
Attraction is itself the positing in the first place of a one distinct from other
ones; these are only the immediate ones that are to preserve themselves
through repulsion; through their posited negation, however, what proceeds
is the one of attraction which is therefore determined as the mediated one,
the one posited as one. The first ones, as immediate, do not in their ideality
return into themselves, but have this ideality each in another.
The one one is, however, ideality that has been realized, posited in the
one; it attracts through the mediation of repulsion; it contains in itself this
mediation as its determination. It thus does not swallow the attracted ones
within it as into one point, that is, does not sublate them abstractly. Since
it contains repulsion in its determination, the latter equally preserves the
ones as many within it; by its attracting, it musters, so to speak, something
before it, gains an area or a filling. Thus there is in it the unity of repulsion
and attraction in general.

c. The connection of repulsion and attraction

The difference of the one and the many has determined itself as a difference
of their mutual reference connecting them which breaks down into two,
repulsion and attraction, each of which stands at first outside the other
on its own, in such a way that the two are essentially joined together
nevertheless. Their still indeterminate unity must be brought out in greater
detail.
As the fundamental determination of the one, repulsion appears first, and
it appears as immediate, like its ones which are indeed generated by it and
yet are at the same time posited as immediate, and it is therefore indifferent
to the attraction which is added to it externally as thus presupposed. Rather,
attraction is not presupposed by repulsion: it is not supposed to have any
part in the positing and in the being of the latter, that is, as if repulsion
were not, already in it, the negation of itself, or the ones were not already
negated in it. In this way, we have repulsion in abstraction, by itself, and
attraction likewise holds out to the ones, as each an existent, the side of an
immediate existence which comes to them by itself as an other.
If we take mere repulsion in this way, for itself, it is then the dispersion
of the many ones in indeterminacy, outside the sphere of repulsion itself;
for repulsion is the negating of the connection of the many to one another;
lack of connection is their determination when abstractly taken. But repul-
sion is not just the void; the ones, although unconnected, do not repel
what constitutes their determination, do not exclude it. Although nega-
tive, repulsion is nonetheless essentially connection; the mutual repulsion
and flight is not a liberation from what is repelled and fled from; that which
is excluded still stands in connection with what is excluded from it. But this
moment of connection is attraction, which is thus within repulsion itself;
it is the negating of that abstract repulsion by which the ones would each
be an existent referring only to itself without mutual exclusion.
But in starting with the repulsion of the ones as immediately present
there, and with attraction consequently also posited as intruding on them
externally, the two, repulsion and attraction, are held apart as diverse
determinations despite their inseparability. But it has been established that
it is not just repulsion which is presupposed by attraction, but that there
equally is present also a reverse connection of repulsion to attraction, and
that repulsion no less has attraction for its presupposition.Being-for-itself
As thus determined, they are inseparable, and at the same time each
is determined as an ought and a limitation with respect to the other.
Their ought is their abstract determinateness as each an existent in itself –
a determinateness, however, which is thereby directed beyond itself and
refers to the other. And so, through the mediation of the other, each is as
other; their self-subsistence consists in their being mutually posited in this
mediation as an other determining. – Thus, repulsion is the positing of the
many; attraction the positing of the one; this latter is equally the negation
of the many and the former the negation of the ideality of such a many in
the one; so that attraction too is attraction only through the mediation of
repulsion, just as repulsion is repulsion through the mediation of attraction.
In all this, however, the mediation of each with itself through the other is
in fact negated; each of the two determinations is its own self-mediation.
This will result from a closer examination of the two determinations and
will bring us back to the unity of their concept.
In the first place, that each presupposes itself, that in its presupposition
each refers only to itself, this is already present in the way the still relative
repulsion and attraction behave at first.
Relative repulsion is the mutual repulsion of many ones which are already
at hand, supposedly immediately given. But that there be many ones, this is
repulsion itself; any presupposition that it would have is only its own posit-
ing. Moreover, the determination of the being that would accrue to the ones
apart from their being posited – whereby they would already be – belongs
likewise to repulsion. Repelling is that through which the ones manifest
themselves and maintain themselves as ones; through which they are as
such. Their being is their repulsion itself, which is thus not some relative
existence against another other but relates itself throughout only to itself.
Attraction is the positing of the one as such, of the real one, with respect
to which the existence of the many is determined as only a vanishing
idealization. Attraction thus directly presupposes itself; it presupposes itself
in the determination namely, of the many ones to be an idealization, the
same ones which are otherwise supposed to have existence for themselves
and to repel others, including therefore any other that attracts. Against this
determination of repulsion, the ones do not attain ideality only through
the relation to attraction; on the contrary, the ideality is presupposed: it
is the ideality of the ones as an existent in itself, inasmuch as they, as ones
(including the one conceived as attracting), are not distinguished from one
another but are one and the same.
This self-presupposing of the two determinations, each for itself,
implies further that each contains within itself the other as moment.
Self-presupposing in general is the positing of oneself in a one as the neg-
ative of oneself (repulsion), and what is presupposed in this positing is
the same as that which presupposes (attraction). That each is in itself only
a moment, this is the transition of each from itself into the other, the
negation of itself in the other and the positing of itself as the other of
itself. The one, as such, is thus a coming-out-of-itself; is itself only the
positing of itself as its other, as the many. And the many, for its part, is
only the falling back upon itself and the positing of itself as its other, as a
one, and is in this equally only the connecting of itself to itself, each con-
tinuing itself in its other. Therefore, the coming-out-of-itself (repulsion)
and the self-positing-as-one (attraction) are already inherently present as
undivided. But in the repulsion and attraction which are relative, that is,
which presuppose immediate, determinedly existent ones, it is posited that
the two are each, within it, this negation of itself, and consequently also the
continuity of itself in its other. The repulsion of the determinedly existent
ones is the self-preservation of the one through the mutual holding off of
the others, so that (1) the other ones are negated in it (this is the side of
its existence or of its being-for-another and is therefore attraction as the
ideality of the ones); and (2) the one is in itself, without reference to the
others (however, not only has the in-itself in general long since passed over
into being-for-itself; the one in itself, according to its determination, is
the coming to be of many). – The attraction of the existent ones is their
ideality and the positing of the one, and in this, as both the negating and
the producing of the one, attraction sublates itself, and as a positing within
it of the one, is the negative of itself: it is repulsion.
With this, the development of being-for-itself is completed and has
attained its result. In connecting itself to itself infinitely, that is, as the
posited negation of negation, the one is the mediation by which it repels
itself as its absolute (that is, abstract) otherness (the many) from itself, and
in thus negatively connecting itself to this, its non-being, it sublates it and is
in it precisely only the connection to itself. The one is only this becoming
in which the determination “it begins,” that is, its being posited as an
immediate existent, and equally that, as result, it has restored itself as the
one, that is, the equally immediate and exclusive one, have vanished; the
process which it is, posits and contains it from all sides only as something
sublated. The sublation, determined at first only as a relative sublating of
the connection to another existent, a connection which is therefore itself not
an indifferent repulsion and attraction, equally proves itself to pass over into
the infinite connection of mediation through the negation of the external
connection of immediate and determinate existents, and to have for result
precisely that becoming which, in the instability of its moments, is the
collapse, or rather the going-together-with-itself, into simple immediacy.
This being, according to the determination which it has now acquired, is
quantity.

If we briefly review the moments of this transition of quality into quantity,
we find that the qualitative has being and immediacy for its fundamental
determination, and the limit and the determinateness are in this immediacy
so identical with the being of something, that the something itself vanishes
along with its alteration; as thus posited, it is determined as finite. Because
of the immediacy of this unity in which the distinction has disappeared,
although it is implicitly present in the unity of being and nothing, the
distinction falls outside that unity as otherness in general. This reference to
the other contradicts the immediacy in which qualitative determinateness
is self-reference. This otherness is sublated in the infinity of the being-for-
itself, the being-for-itself that has realized the distinction implicitly present
in it in the negation of negation: has realized it as the one and the many
and as their connecting references, and has also elevated the qualitative to
true unity, that is, a unity which is no longer immediate but posited as
accordant with itself.

This unity is, therefore,
(a) being, only as affirmative, that is, immediacy
self-mediated through the negation of negation:
being is posited as a unity permeating
its determinacies, limits, etc.,
which are posited in it as sublated;
(b) existence:
in this determination it is negation
or determinateness as moment of
the affirmative being;
yet this determinateness is no longer
immediate but reflected into itself,
refers not to another but to itself;
(c) absolutely-determined-being,
(d) absolute in-itselfness,
(e) the one;
(f) otherness as such is itself being-for-itself;
(g) being-for-itself:
as that being which persists across the determinateness
and in which the one and even the being-determined-in-itself are posited as sublated.
The one is simultaneously determined as having gone beyond itself and as unity;
the one, the absolutely determined limit,
is consequently posited as a limit which is none,
a limit which is in being but is indifferent to it.
